\documentclass[11pt]{amsart}
\usepackage[T1]{fontenc}
\usepackage{lmodern,amsmath,amssymb,amsthm,mathtools,microtype,booktabs,array,longtable}
\usepackage[margin=1.06in]{geometry}
\usepackage{tikz}
\usepackage{xr-hyper}
\usepackage[hidelinks]{hyperref}
\hypersetup{pdftitle={Finite-Use Discrimination Geometry of Ordered Quantum Measurements},pdfauthor={Qian Qi},pdfsubject={General Theta Foundations I, Revision 96}}
\emergencystretch=2em
\numberwithin{equation}{section}
\newtheorem{theorem}{Theorem}[section]
\newtheorem{lemma}[theorem]{Lemma}
\newtheorem{proposition}[theorem]{Proposition}
\newtheorem{corollary}[theorem]{Corollary}
\theoremstyle{definition}\newtheorem{definition}[theorem]{Definition}
\theoremstyle{remark}\newtheorem{remark}[theorem]{Remark}
\newcommand{\Q}{\mathbb Q}\newcommand{\R}{\mathbb R}\newcommand{\C}{\mathbb C}\newcommand{\N}{\mathbb N}\newcommand{\Z}{\mathbb Z}\newcommand{\E}{\mathbb E}
\newcommand{\ip}[2]{\langle #1,#2\rangle}\newcommand{\norm}[1]{\lVert #1\rVert}
\newcommand{\epsc}{\varepsilon_{\mathrm c}}
\DeclareMathOperator{\osc}{osc}\DeclareMathOperator{\spanop}{span}\DeclareMathOperator{\End}{End}\DeclareMathOperator{\diag}{diag}\DeclareMathOperator{\tr}{tr}\DeclareMathOperator{\rad}{rad}\DeclareMathOperator{\TV}{TV}
\title[Finite-use discrimination geometry]{Finite-Use Discrimination Geometry of Ordered Quantum Measurements}
\author{Qian Qi}
\date{6 October 2026}
\subjclass[2020]{Primary 81P15, 81P47; Secondary 94A17, 81P70}
\externaldocument{build/xref-supplement}[BINARY_SUPPLEMENT.pdf]
\begin{document}
\raggedbottom
\begin{abstract}
We study finite decision experiments on ordered quantum measurements
when each fresh acquisition is conditionally independent of the retained
receiver.  For arbitrary finite two-call experiments we give an attained
common-barycenter variational principle and a Hermitian-majorant dual,
including specified retained-register dimensions.  A Heisenberg-picture
postponement lemma identifies the value without a receiver message.
With the complete old Schmidt support retained, messages are useless at
every prescribed pure first acquisition exactly when each labelwise
continuation value is concave.  If this condition fails, a two-outcome
receiver instrument after one device label witnesses an advantage at
some first acquisition.  For a weighted shared-basis experiment with a
specified biased prior, we evaluate every prescribed-initial-spectrum
value through the least rank-capped majorant and a secular equation.
An instrument with at most $d$ outcomes attains it without changing the
initial acquisition.  We obtain its exact message-gain criterion,
saturation and rigidity.  For equal priors and a two-dimensional old
receiver we also give the full initial-spectrum formula in every input
dimension, including the complete qubit case.  The arbitrary-experiment
characterization and the shared-basis closed formulas have different
quantifiers.  The previously established local tangent-tube, profile
and hard-budget geometry is retained in the linked supplement.
\end{abstract}
\maketitle
\enlargethispage{3pt}
\section{Introduction}
\label{sec:introduction96}
A quantum measurement can produce only a classical label while leaving
information in a reference prepared by the experimenter.  Across two
uses, retaining this receiver, sending its classical measurement result
to a fresh source, and coherently feeding it into a later acquisition
are different operations.  This article studies their exact decision
values in a fixed operational interface.

An ordered measurement is a tuple $E=(E_y)_y$ of positive matrices with
$\sum_yE_y=I$.  Its unknown interface is
$\mathcal M_E(X)=\sum_y\tr(E_yX)|y\rangle\langle y|$;
no residual device output is accessible.  A reset cut requires the
\emph{entire} new input--reference system to be conditionally independent
of the old receiver.  The old and fresh references may be measured
jointly only after the second call.  All controls are common to the
hypotheses, and all formal histories carry a defined rule.
We use unhalved trace norms; a decision-probability perturbation carries
the corresponding factor $1/2$ for a trace-zero difference.

\subsection{General experiments and the value of a message}
For arbitrary finite priors and rewards, the branch normal form is
$\sum_h A_{yh}=\rho$, with a fresh density $b_{yh}$ on each branch.
The trace of $A_{yh}$ is a Gram-decomposition weight, not a physical
history probability.  Optimizing the fresh state and the final POVM
gives a continuation function $g_y$; receiver messages replace it by
its upper concave envelope $c_y$.  Theorems~\ref{thm:resetvariational91}
and~\ref{thm:resetdual91} give attained primal and dual values, with
finite instruments and universal Hermitian majorants.  Their
rank-resolved forms specify which complete quantum registers are
counted at each cut.

Lemma~\ref{lem:postponement96} proves the no-message reduction directly
in the Heisenberg picture.  Theorem~\ref{thm:messagecriterion96} makes
its uniform consequence exact: with the full old Schmidt support
retained, receiver messages have no value at every prescribed pure
first acquisition if and only if every fresh-rank-constrained
continuation is concave.  Failure is witnessed by two outcomes after
one first device label at some initial acquisition.  This assertion
is not a two-message bound at every preassigned initial state.
Nor may postponement bypass an old-dimension bound smaller than the
system it retains.

\subsection{A solved finite family}
We next use a weighted exact second-moment family of ordered bases.
An alternative basis is selected once and reused at both calls;
independent redraws give a different experiment with no second-moment
signal.  Theorem~\ref{thm:operationalhierarchy90} evaluates the three
unrestricted endpoint classes, and Theorem~\ref{thm:rankspectrum92}
resolves both retained-register dimensions.

For the specified biased prior and a prescribed pure initial Gram
$\rho$, Theorem~\ref{thm:spectralvalue94} gives
\[
 P_{k,\ell}^{\rm b}(\rho)
  =\frac{d+1}{2(d+1)-t^2}
       +\frac{t^2}{2[2(d+1)-t^2]}
             \chi\bigl(q^{(\min(k,\ell))}(\lambda(\rho))\bigr).
\]
The boundary-safe definitions of $\chi$ and the unique least rank-capped
majorant $q^{(r)}$, including ties, are given before that theorem in
Section~\ref{sec:spectralvalue94}.  A finite instrument on the fixed
initial receiver attains the formula.  No-message preparation instead
uses the leading truncated spectrum.  Thus at $t>0$ the message has
strict value exactly for $2\le r<\operatorname{rank}\rho$.
Theorem~\ref{thm:initialspectra93} and the subsequent rigidity estimates
describe equality and quantitative loss.  Proposition~\ref{prop:qutritwitness96}
works out a complete three-dimensional example, with exact algebraic
values and rational secular intervals.

The prior matters.  For equal priors and old receiver dimension two,
Theorem~\ref{thm:equalinitial95} gives a different initial-spectrum
formula and different optimal fresh Schmidt weights.  Its qubit
specialization covers every pure first acquisition without a receiver
message.  The full proof includes the optimizing fresh weights, the singular
boundary and finite attainment on the prescribed receiver.

\subsection{Relation to local geometry}
The same ordered-measurement interface also supports local fixed-tube
results in the linked supplement.  Those results classify support and
covariance directions, preserve the finite tube remainders, and give
allocation orders involving $N=\sum_jn_j$, $V=\sum_jn_j^2$, and hard
pathwise $Q$.  They are not used to prove the exact finite-game theorems.
Conversely the spectral formulas do not make their fixed-object
constants uniform over changing support strata.  The local theory is
retained as mathematics, not discarded; it is separated from the
finite-experiment proof route.  The independent structural article and
the wider analytic Theta programme supply no hidden hypotheses here.

\section{A variational principle for two-call reset experiments}
\label{sec:resetvariational91}
The fresh-preparation cut admits an exact finite-dimensional description
of optimal decision values.  No symmetry, second-moment identity, or
special choice of prior is needed in this section.  The description is
by a concave envelope over density matrices, not over pure states alone.
Its finite realization bounds the receiver needed to attain a value;
it does not identify reset width with a prescribed memory dimension.

\subsection{Experiments and strategy classes}
Let $I_\ell=\C^{d_\ell}$ and let $Y_\ell$ be a finite classical alphabet,
$\ell=1,2$.  A finite experiment consists of a probability law
$(p_\theta)$, two ordered measurements $E^{\theta,\ell}$ on $I_\ell$,
and rewards $r_{j\theta}\in[0,1]$ for decisions $1\le j\le m$.
The index $\theta$ is sampled once.  Conditional on it the two slots are
the memoryless measurement channels $\mathcal M_{E^{\theta,\ell}}$.
Repeating one unknown measurement is the case
$E^{\theta,1}=E^{\theta,2}$.  Identification and binary decisions on
subfamilies are obtained by taking indicator rewards.

\begin{definition}[Two-call classes]\label{def:classes91}
An unrestricted strategy prepares a state on $I_1\otimes R$, observes
the first label $y$, applies a common trace-preserving map from $R$ to
$I_2\otimes R'$, and, after observing $z$, measures $R'$ to decide $j$.
A unit-reset strategy instead applies a normal instrument
$(\mathcal K_{yh})_h$ to $R$, retains its old quantum output $O_{yh}$,
and prepares a state $\sigma_{yh}$ on the complete fresh system
$I_2\otimes R_{2,yh}$.  Conditional on $(y,h)$ the preparation is
exactly the tensor product of the old output with $\sigma_{yh}$.
No old quantum register enters $I_2$ or its fresh reference.  After the
second label, the two receivers may be measured jointly.  A classically
adaptive strategy is a unit-reset strategy that measures the entire old
receiver and retains only a classical record before fresh preparation.

All rules are independent of $\theta$ and defined on every formal
history, including histories null under a hypothesis.  Receiver spaces
have no dimension bound.  Instruments may have countably many outcomes;
all classical records are retained.  Finite or countable public
randomization and the
finite-dimensional trace-norm closure of the resulting output-dephased
Choi tester tuples are included in each class.  General public seed laws
are included through this closure: their finitely many tester coordinates
are limits of finite convex combinations.  Early stopping is
implemented by a fixed, discarded fresh call and a decision independent
of its label; it is not a zero-length component of a prescribed profile.
Write $P_{\rm all},P_{\rm reset},P_{\rm ca}$ for the suprema of expected
reward in these classes.
\end{definition}
This reset cut is the two-call instance of Definition~\ref{def:reset88}.
The first class in the ordering
$P_{\rm ca}\le P_{\rm reset}\le P_{\rm all}$ agrees with complete
measure-and-reprepare adaptation, not with a constraint of separability
on all final reset effects.  In the order $I_1,Y_1,I_2,Y_2$, a
classically adaptive effect is a sum of positive products across
$I_1Y_1\mid I_2Y_2$ \cite[Definition~23]{OZNPQ89}.
The positive separable cone is closed in these finite dimensions:
its trace-one slice is the compact convex hull of product states, and
trace controls passage to the cone.  Only this necessary condition,
not its sufficiency for classical adaptive normalization, is used below.

\begin{lemma}[Countable records]\label{lem:countable91}
Let $(X_h)$ be positive trace-class operators with
$\sum_h\tr X_h<\infty$.  Their partial sums converge in trace norm,
and their recorded direct sum lies in the trace-class $\ell^1$ sum.
For effects $0\preceq F_h\preceq I$, the corresponding probability
series converges absolutely.  Replacing all but finitely many branches
of a normal instrument by one remainder branch preserves normalization;
changing the continuation on that branch changes any reward in $[0,1]$
by at most its probability.
\end{lemma}
\begin{proof}
The norm of a positive tail equals its trace.  This proves the first
assertions and bounds the probability tail by $\sum_{h>n}\tr X_h$.
For an instrument, the sum of the omitted completely positive maps is
the remainder map; its trace is precisely the omitted outcome
probability.  Apply the bound to each member of the finite experiment
and then average.  The maximum of its finitely many tail probabilities
tends to zero.  No conditional state is divided by a history probability.
\end{proof}

\subsection{A complete branch normal form}
Fix bases on $I_1,I_2$.  A coefficient map $C:I_1\to R$ represents
$|C\rangle=\sum_i|i\rangle\otimes C|i\rangle$; its squared norm is
$\tr C^*C$.  Measuring its input gives the unnormalized receiver state
$CE_y^{\mathsf T}C^*$.

\begin{lemma}[Unit-reset normal form and converse]\label{lem:resetnormal91}
Every implemented unit-reset strategy has a purification, still respecting the reset
cut, from which its original decision statistics are recoverable by
common receiver channels.  For this purification there are rectangular
maps
\[
 C_0:I_1\longrightarrow R_0,\quad
 C_{yh}:I_1\longrightarrow O_{yh},\quad
 D_{yh}:I_2\longrightarrow R_{2,yh}
\]
with, on the input copies,
\begin{equation}\label{eq:resetgrams91}
 \rho=C_0^*C_0,\quad \tr\rho=1,\quad
 A_{yh}=C_{yh}^*C_{yh},\quad \sum_hA_{yh}=\rho,\quad
 b_{yh}=D_{yh}^*D_{yh},\quad\tr b_{yh}=1.
\end{equation}
The receiver block on the formal record $(y,h,z)$ is
\begin{equation}\label{eq:resetblock91}
 (C_{yh}E_y^{\theta,1\,\mathsf T}C_{yh}^*)
 \otimes(D_{yh}E_z^{\theta,2\,\mathsf T}D_{yh}^*).
\end{equation}
The sum in \eqref{eq:resetgrams91} converges in trace norm.
Conversely, every finite collection satisfying \eqref{eq:resetgrams91}
is realizable, up to common receiver isometries, by a unit-reset strategy
with old receiver dimension at most $d_1$ and fresh reference dimension
at most $d_2$.  Zero-probability histories require no extra assumption.
The tester closure in Definition~\ref{def:classes91} is used by continuity
of decision values; no physical dilation is assumed for an abstract
limit before its value has been realized by the finite optimization.
\end{lemma}
\begin{proof}
Purify the first preparation, retaining the purifying system in $R_0$.
The Schmidt support has dimension at most $d_1$.  For each $y$, dilate
each branch map of the intervening instrument to a single map
$V_{yh}:R_0\to O_{yh}$, retaining its environment as old receiver
memory.  On the relevant support,
$\sum_hV_{yh}^*V_{yh}=I$ in the strong sense.  Thus
$C_{yh}=V_{yh}C_0$ has
$\sum_h C_{yh}^*C_{yh}=C_0^*C_0$.
The positive tails converge in trace norm by Lemma~\ref{lem:countable91}.
This argument uses only the finite Schmidt support even when the
original receiver or a dilation environment is infinite-dimensional.
Tracing those environments recovers the original instrument.

Purify the entire fresh preparation separately.  Its coefficient map
is $D_{yh}$ with squared Hilbert--Schmidt norm one.  The tensor-product
preparation rule, not a statement about a probe marginal, gives
\eqref{eq:resetblock91}.  Common pre-query controls are absorbed into
these preparations and common post-query controls into the final
receiver measurement.  Fresh mixing may be purified inside the fresh
reference.  A retained public seed may be put in the classical record;
randomization before the first call is a mixture of such strategies and
cannot increase the supremum of a linear reward.  More explicitly,
for seed law $q_s$ and first coefficient maps $C_{0,s}$, take
$C_0=\bigoplus_s\sqrt{q_s}C_{0,s}$ and measure the seed register
as part of the first receiver instrument.  The branch maps are
$C_{y,(s,h)}=\sqrt{q_s}V_{ysh}C_{0,s}$, and their Gram sum is
$\sum_sq_sC_{0,s}^*C_{0,s}=C_0^*C_0$ for every $y$.
The fresh rule retains its dependence on $(s,y,h)$, entirely on the
fresh side of the cut.  Thus the original seeded decision statistics
are recovered.  Countable records and tester limits follow from the
preceding lemma and continuity.

For the converse let $r=\operatorname{rank}\rho$ and choose
$C_0:I_1\to\C^r$ with $C_0^*C_0=\rho$.  Its restriction to
$\operatorname{supp}\rho$ is invertible onto $\C^r$.  Let $C_0^+$
be the inverse on this support, and set
\[
 C_{yh}=\sqrt{A_{yh}},\qquad
 V_{yh}=\sqrt{A_{yh}}C_0^+,\qquad D_{yh}=\sqrt{b_{yh}}.
\]
Positivity and $\sum_hA_{yh}=\rho$ imply
$\operatorname{supp}A_{yh}\subseteq\operatorname{supp}\rho$.
Hence $V_{yh}C_0=\sqrt{A_{yh}}$ and
$\sum_hV_{yh}^*V_{yh}=I_r$.  These are Kraus operators of a
trace-preserving recorded instrument with outputs in $\C^{d_1}$.
The fresh vector $|\sqrt{b_{yh}}\rangle$ is normalized and independent
of that output.  Polar decomposition identifies any other rectangular
maps of the same Gram matrices with these maps on their supports;
receiver isometries may be absorbed into the final effects.
Zero Gram matrices give zero branch maps.  All other null histories are
handled by the same unnormalized formulas, including projective devices.
\end{proof}

\subsection{The compact primal problem}
Let $\mathcal D_d$ denote the density matrices on $\C^d$, and put
\begin{equation}\label{eq:gameblocks91}
 G_{j,yz}=\sum_\theta p_\theta r_{j\theta}
     (E_y^{\theta,1})^{\mathsf T}\otimes(E_z^{\theta,2})^{\mathsf T}.
\end{equation}
These are fixed positive matrices on the two input copies.  For
$A\succeq0$ and $b\in\mathcal D_{d_2}$ define
\begin{equation}\label{eq:leafvalue91}
 \Phi_y(A,b)=\sum_z\max_{\substack{F_j\succeq0\\\sum_jF_j=I}}
 \sum_j\tr\bigl[F_j(\sqrt A\otimes\sqrt b)
            G_{j,yz}(\sqrt A\otimes\sqrt b)\bigr].
\end{equation}
The POVM is chosen separately at every recorded $z$.  It acts on
$\C^{d_1}\otimes\C^{d_2}$.  For $a\in\mathcal D_{d_1}$ let
$g_y(a)=\max_{b\in\mathcal D_{d_2}}\Phi_y(a,b)$ and let
\begin{equation}\label{eq:roof91}
 c_y(\rho)=\sup\left\{\sum_h\lambda_h g_y(a_h):
 \lambda_h\ge0,\ \sum_h\lambda_h=1,\
 \sum_h\lambda_ha_h=\rho\right\}.
\end{equation}
The envelope is taken over \emph{all} density matrices $a_h$.  Replacing
them by rank-one states would change the resource problem.  The
construction of concave envelopes is standard \cite{Uhlmann91}; the
operational content here is the reset realization and the common
barycenter constraint across all first labels.

\begin{theorem}[Exact reset variational principle]\label{thm:resetvariational91}
For every finite experiment of Definition~\ref{def:classes91},
\begin{equation}\label{eq:resetvariational91}
 P_{\rm reset}=\max_{\rho\in\mathcal D_{d_1}}\sum_yc_y(\rho)
 =\max_{\substack{A_{yh}\succeq0,\ \sum_hA_{yh}=\rho\\
                    \rho\in\mathcal D_{d_1},\ b_{yh}\in\mathcal D_{d_2}}}
                    \sum_{y,h}\Phi_y(A_{yh},b_{yh}).
\end{equation}
All maxima are attained.  For a fixed $\rho$ of rank $r$, at most $r^2$
nonzero instrument outcomes are required after each first label.
An optimal strategy exists with old receiver dimension at most $d_1$,
fresh reference dimension at most $d_2$, and a final quantum receiver
of dimension at most $d_1d_2$, separately from the finite classical
record.  The dimension bounds concern the retained registers at the recording
cuts, not temporary workspaces for implementing an instrument.
These assertions include arbitrary priors and arbitrary fixed
pairs as particular decision problems; they assert an exact optimum,
not a strict advantage for every such problem.
\end{theorem}
\begin{proof}
For fixed $y$, $\Phi_y$ is continuous, and is positively homogeneous
in $A$.  Continuity follows from continuity of matrix square roots and
maximization of a continuous function over a compact POVM set.
Maximizing also over compact $\mathcal D_{d_2}$ makes $g_y$ continuous.
The graph $\{(a,g_y(a)):a\in\mathcal D_{d_1}\}$ is compact in a
real space of affine dimension at most $d_1^2$.  Its convex hull is
compact.  The largest height above $\rho$ is therefore attained and
is exactly $c_y(\rho)$, also allowing countable decompositions.
It is an upper semicontinuous concave function of $\rho$.
Consequently $\sum_yc_y$ attains a maximum.

Carath\'eodory first supplies at most $r^2+1$ atoms at a maximizing
height over a rank-$r$ matrix: each atom is supported on
$\operatorname{supp}\rho$, since positive matrices cannot cancel on its
orthogonal complement.  If more than $r^2$ weights are nonzero, the
Hermitian atoms are linearly dependent.  Choose nonzero real $v_h$ with
$\sum_hv_ha_h=0$; tracing gives $\sum_hv_h=0$.
Small perturbations $\lambda_h\pm\epsilon v_h$ preserve the barycenter
and positivity.  Optimality forces $\sum_hv_hg_y(a_h)=0$, since otherwise
one sign raises the height.  Move in this direction until a weight
vanishes.  The value is unchanged.  Iteration gives at most $r^2$ atoms.
For each atom the maximum over $b$ and the leaf POVMs is attained.

By Lemma~\ref{lem:resetnormal91}, the value of any purified competitor
is bounded by the right side of \eqref{eq:resetvariational91}, with
$A_{yh}=\lambda_{yh}a_{yh}$.  Countable records follow by the positive
tail bound; public mixtures and norm limits preserve the upper.
Conversely, the lemma realizes every finite maximizing collection,
and its optimal leaf POVMs realize every maximum in
\eqref{eq:leafvalue91}.  This proves equality and the dimension bounds.
The final quantum dimension excludes the separately recorded labels.
\end{proof}
For a binary equal-prior fixed pair, its optimal unhalved final trace
separation is $4P_{\rm reset}-2$.  Thus the theorem is also a general
variational formula for that operational distance.  The primal has
finitely many bounded matrix variables, but is not asserted to be a
convex semidefinite program or an efficient synthesis algorithm.

\subsection{Dual certificates and the role of receiver feedback}
\begin{theorem}[Duality and contact conditions]\label{thm:resetdual91}
For the same experiment,
\begin{equation}\label{eq:resetdual91}
 P_{\rm reset}=\min_{(H_y)}\lambda_{\max}\left(\sum_yH_y\right),
 \qquad \tr(H_ya)\ge g_y(a)\quad(a\in\mathcal D_{d_1}),
\end{equation}
where $H_y$ are Hermitian $d_1\times d_1$ matrices.  The minimum is
attained.  It is equivalent to minimizing $\lambda$ subject to the
same majorizations and $\sum_yH_y=\lambda I$.
For any optimal primal and dual, every atom of positive weight is a
contact point, $\tr(H_ya_{yh})=g_y(a_{yh})$, and $\rho$ is supported
on the maximal eigenspace of $\sum_yH_y$.  Conversely these conditions,
with maximizing fresh states and leaf POVMs, certify optimality.
\end{theorem}
\begin{proof}
Every majorization of $g_y$ majorizes its concave envelope.  Therefore
$\sum_yc_y(\rho)\le\tr(\rho\sum_yH_y)\le
\lambda_{\max}(\sum_yH_y)$, which proves weak duality.
For completeness, the strong-duality step is finite-dimensional.
Extend $c_y$ homogeneously to positive matrices by
$\widehat c_y(A)=\tr A\,c_y(A/\tr A)$ and $\widehat c_y(0)=0$.
Its hypograph
\[
 \mathcal K_y=\{(A,u):A\succeq0,\ u\le\widehat c_y(A)\}
\]
is a closed convex cone.  Closedness at $A=0$ follows from boundedness
of $g_y$; elsewhere it follows from the compact-graph construction.
The conic problem maximizing $\sum_yu_y$, with $(A_y,u_y)\in\mathcal K_y$,
$A_y=\rho$ and $\tr\rho=1$, is precisely the primal problem.
The point $A_y=\rho=I/d_1$, $u_y=-1$ is strictly interior to all
hypographs, since rewards are nonnegative and $A_y$ is positive definite.
Its value is finite, in fact at most one by the operational realization.
The finite-dimensional separating-hyperplane form of conic Slater
duality thus gives an attained dual with no gap
\cite[Sections~5.2.3 and~5.9.1]{BV91}.

Its multipliers can be read directly: the supremum of
$u-\tr(H_yA)$ on $\mathcal K_y$ is zero exactly when
$\tr(H_ya)\ge g_y(a)$ for every density matrix.  Taking $\rho$ as a
free Hermitian variable under its trace constraint imposes
$\sum_yH_y=\lambda I$.  This gives the stated dual.  Conversely, for
majorants with sum $S$, add the positive matrix
$\lambda_{\max}(S)I-S$ to one $H_y$; all majorizations remain valid
and the sum becomes scalar.  The two forms have the same minimum.
Finally, the weak-duality difference is a sum of nonnegative atom
majorization defects and a nonnegative spectral defect.  It vanishes
exactly under the stated contact and support conditions.
\end{proof}
The majorizations in \eqref{eq:resetdual91} are universal inequalities
on density matrices.  A finite sampling of them is not a dual certificate.

\begin{corollary}[Quantitative contact certificates]\label{cor:contactdefect91}
Let $(H_y)$ be any feasible all-density majorants in
Theorem~\ref{thm:resetdual91}, put $S_H=\sum_yH_y$ and
$U=\lambda_{\max}(S_H)$, and consider a finite or countable reset
normal form with achieved reward $v$.  For every nonzero branch write
$w_{yh}=\tr A_{yh}$, $a_{yh}=A_{yh}/w_{yh}$, and let $v_{yh}$
be its contribution to the reward, summed over the second labels and
decisions.  Then
\begin{equation}\label{eq:contactdefect91}
\begin{split}
 U-v={}&\tr[\rho(UI-S_H)]\\
 &+\sum_{y,h}w_{yh}\bigl[\tr(H_ya_{yh})-g_y(a_{yh})\bigr]\\
 &+\sum_{y,h}\bigl[w_{yh}g_y(a_{yh})-v_{yh}\bigr].
\end{split}
\end{equation}
Every term is nonnegative.  Consequently, if $U-v\le\varepsilon$,
the total Gram weight of branches with majorization defect at least
$u>0$ is at most $\varepsilon/u$.  If the maximal eigenspace projection
$P$ of $S_H$ has a spectral gap $\gamma>0$ to its complement, then
$\tr[(I-P)\rho]\le\varepsilon/\gamma$.
\end{corollary}
\begin{proof}
The common-barycenter constraints give
$\sum_{y,h}w_{yh}\tr(H_ya_{yh})=\tr(\rho S_H)$.
Add and subtract this quantity and $\sum_{y,h}w_{yh}g_y(a_{yh})$.
The leaf optimization bounds $v_{yh}$ by $w_{yh}g_y(a_{yh})$.
Feasible majorization and $S_H\preceq UI$ prove nonnegativity.
All series converge absolutely: the first-label alphabet is finite,
$\sum_hw_{yh}=1$ for every $y$, and the functions and matrices are
bounded.  Restricting a nonnegative sum proves the first estimate;
$UI-S_H\succeq\gamma(I-P)$ proves the second.
\end{proof}
The weights in this certificate are input Gram weights, not the
hypothesis-dependent probabilities of recorded histories.  With an
attained optimal dual, $U-v$ is precisely the loss in optimal reward.
The formula requires a genuine universal majorant, not sampled tests
of its inequality.

\begin{corollary}[Label feedback versus receiver feedback]\label{cor:feedbackroof91}
If no receiver-instrument outcome is used to choose the second
preparation, but the first device label may be used, the optimum is
\[
 P_{\rm label}=\max_{\rho\in\mathcal D_{d_1}}\sum_yg_y(\rho).
\]
Allowing receiver feedback replaces $g_y$ by $c_y$ in this formula.
In particular, concavity of every $g_y$ is sufficient for
$P_{\rm label}=P_{\rm reset}$.  More generally, a label-only strategy
with a common density $\rho$ is optimal for the full reset class if
it attains the contact and spectral conditions of
Theorem~\ref{thm:resetdual91}.
\end{corollary}
\begin{proof}
With one preparation choice after each $y$, the normal form has
$A_y=\rho$; a trace-preserving receiver operation cannot improve a
subsequent unrestricted joint measurement and can be postponed to it.
Conversely these choices are admissible.  The assertion follows from
the primal formula and the definition of concave envelope.
\end{proof}

\subsection{A geometric criterion for retained receiver advantage}
Let $\mathcal P_{d_1}$ be the rank-one density matrices.  In parallel
with \eqref{eq:roof91}, define
\begin{equation}\label{eq:pureroot91}
 c_y^{\rm cl}(\rho)=\sup\left\{\sum_h\lambda_h g_y(a_h):
 a_h\in\mathcal P_{d_1},\quad
 \lambda_h\ge0,\quad\sum_h\lambda_h a_h=\rho\right\}.
\end{equation}
The weights sum to one by taking traces.  Every density matrix admits
such a decomposition, but the value of this pure-atom envelope need
not equal the all-density envelope $c_y$.

\begin{theorem}[Classical and quantum receiver envelopes]
\label{thm:classicalroof91}
For every finite experiment of Definition~\ref{def:classes91},
\begin{equation}\label{eq:classicalroof91}
 P_{\rm ca}=\max_{\rho\in\mathcal D_{d_1}}\sum_yc_y^{\rm cl}(\rho)
 =\min_{(H_y)}\lambda_{\max}\left(\sum_yH_y\right),
 \qquad \tr(H_ya)\ge g_y(a)\quad(a\in\mathcal P_{d_1}).
\end{equation}
Both optima are attained.  At most
$(\operatorname{rank}\rho)^2$ nonzero first-receiver outcomes per
first device label suffice at an optimal $\rho$.
The passage from complete classical readout to retained quantum
receiver memory replaces the pure-atom envelopes $c_y^{\rm cl}$ by
the all-density envelopes $c_y$.  In particular, the invariant
\begin{equation}\label{eq:receivergap91}
 \Gamma=\max_\rho\sum_yc_y(\rho)
          -\max_\rho\sum_yc_y^{\rm cl}(\rho)
\end{equation}
is exactly the optimal retained-receiver advantage in that experiment.
\end{theorem}
\begin{proof}
After complete first readout, a classical branch has a positive input
Gram effect $A$ and retains no quantum system depending on $\theta$.
Its probability under a first effect is $\tr(AE_y^{\mathsf T})$.
Refine a spectral decomposition $A=\sum_i\lambda_i a_i$ into the
classical record and keep the former second preparation and decision
on each refined branch.  This does not change the score and gives a
rank-one decomposition with the same barycenter.  Optimizing the
second action separately after the refined record can only increase it.

Conversely, for a rank-one atom $a=|u\rangle\langle u|$, the old
receiver in the canonical normal form has
\[
 \sqrt a E_y^{\mathsf T}\sqrt a
 =\langle u,E_y^{\mathsf T}u\rangle a.
\]
Its normalized state is independent of the unknown device.  A joint
POVM with this known factor is equivalent to a POVM on the fresh
reference alone, by taking its expectation against $a$.
Thus $g_y(a)$ is attainable with only a classical first record.
The converse part of Lemma~\ref{lem:resetnormal91}, followed by discarding
this known old state, realizes every finite pure-atom decomposition.
These two observations give the first equality in
\eqref{eq:classicalroof91}.

The graph of $g_y$ on the compact pure-state set is compact.  The
compact-convex-hull and rank-face atom-elimination arguments of
Theorem~\ref{thm:resetvariational91} therefore apply verbatim, giving
attainment and the stated branch count.  Its homogeneous hypograph is
again a closed convex cone, with a strictly feasible point at a
positive definite common barycenter and negative heights.  The proof
of Theorem~\ref{thm:resetdual91} applies with majorization required only
on pure atoms and gives the attained dual.  Majorization on pure atoms
is equivalent to majorization of $c_y^{\rm cl}$ on every density matrix,
not to majorization of $g_y$ on every density matrix.  The last assertion
now follows from the two exact primal formulas.
\end{proof}

\begin{corollary}[Complete separation certificates]\label{cor:gapcertificate91}
There is a strict retained-receiver advantage if and only if there exist
a finite feasible reset collection from \eqref{eq:resetvariational91},
with achieved reward $v$, and Hermitian matrices $H_y$ such that
\[
 \tr(H_ya)\ge g_y(a)\quad(a\in\mathcal P_{d_1}),\qquad
 v>\lambda_{\max}\left(\sum_yH_y\right).
\]
There is no advantage if and only if an attaining classical strategy
has contact and spectral equality with an all-density dual majorant
of Theorem~\ref{thm:resetdual91}.
\end{corollary}
\begin{proof}
The displayed inequalities compare a physically attained reset lower
with a valid classical upper.  If the gap is positive, take an optimal
finite reset collection and an optimal classical dual, whose existence
was proved above.  This proves both directions of the first assertion.
If the gap is zero, an optimal classical strategy also attains the reset
value; complementarity with any attained reset dual gives the second
condition.  Conversely, that condition certifies a common upper and
lower for the two classes, so their values agree.
\end{proof}
These are certificates with finitely many matrix variables and universal
majorization requirements, not finite-sampling tests.  The criterion
applies to any fixed-pair problem as well as to ensemble decisions,
but does not require every experiment to have positive $\Gamma$.
For example, the equal-prior qubit calculation below gives pure-atom
majorants $H_0=H_1=(7/24)I$, whereas a reset strategy attains $5/8$.
The gap is $1/24$.  Its universal pure-atom majorization follows from
the classical upper in Theorem~\ref{thm:equalprior91}, not from the
finite arithmetic regression.

\subsection{The unrestricted causal normalization}
\begin{proposition}[Physical two-call tester blocks]\label{prop:causalnormal91}
In the unnormalized convention
$J^\sharp(E)=\sum_yE_y\otimes|y\rangle\langle y|$, every unrestricted
strategy has positive blocks $T_{j,yz}$ and positive matrices $\Xi_y$ with
\begin{equation}\label{eq:causalnormal91}
 \sum_jT_{j,yz}=\Xi_y,\qquad
 \tr_{I_2}\Xi_y=\rho_1,\qquad\tr\rho_1=1.
\end{equation}
Its probability of decision $j$ on the output block $(y,z)$ is
$\tr[T_{j,yz}(E_y^{\theta,1}\otimes E_z^{\theta,2})]$.
In particular $0\preceq T_{j,yz}\preceq\Xi_y$ and $\tr\Xi_y=1$.
\end{proposition}
\begin{proof}
Purify the first preparation to $|\psi\rangle\in I_1\otimes R$.
For its common intervening channel $\mathcal A_y:R\to I_2\otimes R'$
put
\[
 X_y=(\operatorname{Id}\otimes\mathcal A_y)(|\psi\rangle\langle\psi|),
 \quad\Xi_y=\tr_{R'}X_y,\quad\rho_1=\tr_R|\psi\rangle\langle\psi|.
\]
For the final POVM $(F_{j,yz})_j$ on $R'$ define
\[
 T_{j,yz}=\tr_{R'}[(I\otimes\sqrt{F_{j,yz}})X_y
                              (I\otimes\sqrt{F_{j,yz}})].
\]
Positivity is immediate.  Cyclicity in the traced factor and
$\sum_jF_{j,yz}=I$ give the sum; trace preservation of $\mathcal A_y$
gives the marginal.  Expanding the Born rule gives the contraction.
Here $\rho_1$ is the physical input marginal, whereas the coefficient
Gram matrix $\rho$ in \eqref{eq:resetgrams91} is its transpose.
These two conventions must not be interchanged.  Every complementary
effect is positive by the same construction.  No factor $d_1d_2$
occurs, because $J^\sharp$ uses the unnormalized maximally entangled
vector.  Countable records and limits follow by continuity.
\end{proof}

\section{Receiver messages and concavification}
\label{sec:messages96}
The distinction between a device label and a receiver message is
operational, not merely terminological.  A device label may select a
fresh source even when an intermediate receiver result may not.  We
first give the exact postponement statement, then characterize when
this additional communication has value in an arbitrary finite
experiment.  The old receiver is large enough to retain the complete
initial Schmidt support throughout this section.  Only the fresh
reference is dimension constrained.

\subsection{Postponement without a receiver message}
Let $R$ be the initial receiver, $I_2$ the second input, and $F_y$ the
fresh reference after the first device label $y$.  The prescribed fresh
state $\sigma_y$ on $I_2\otimes F_y$ is independent of a receiver outcome
$h$.  A normal instrument $\mathcal K_{h|y}:\mathcal T(R)\to
\mathcal T(O_{yh})$ acts on the old receiver.  Write
$M_{j|yhz}$ for the final decision POVM on $O_{yh}\otimes F_y$.
The record $(y,h,z)$ is classical; retaining $h$ for the final readout
is permitted even though it is not sent to the source.

\begin{lemma}[No-message postponement]\label{lem:postponement96}
For every such instrument and readout, the operators
\begin{equation}\label{eq:heisenberg96}
 G_{j|yz}=\sum_h
       (\mathcal K_{h|y}^{*}\otimes\operatorname{Id}_{F_y})(M_{j|yhz})
\end{equation}
form a POVM on $R\otimes F_y$.  Replacing the instrument and its record
by this final POVM preserves every decision probability, under each
hypothesis separately, for every input state and every second device.
The statement includes countable normal instruments and their
finite-dimensional tester closure.  It preserves any fresh-reference
dimension bound.  It preserves an old-receiver bound $k$ whenever
$\dim R\le k$, or after compression to a common input support of
dimension at most $k$.
\end{lemma}
\begin{proof}
Every summand in \eqref{eq:heisenberg96} is positive.  Since the sum of
the instrument maps is trace preserving,
\[
 \sum_jG_{j|yz}=\sum_h\mathcal K_{h|y}^{*}(I_{O_{yh}})
                                      \otimes I_{F_y}=I_R\otimes I_{F_y}.
\]
Let $X_{\theta,y}$ be an arbitrary unnormalized first receiver block
and let $Z_{\theta,y,z}$ be the fresh reference block obtained by
applying the second device to $\sigma_y$.  Conditional independence
at the preparation cut gives
\[
 \begin{split}
 &\sum_h\tr\bigl[M_{j|yhz}
        (\mathcal K_{h|y}(X_{\theta,y})\otimes Z_{\theta,y,z})\bigr]\\
 &\hspace{35mm}=\tr[G_{j|yz}(X_{\theta,y}\otimes Z_{\theta,y,z})].
 \end{split}
\]
This is the adjoint identity, with no division by a history
probability.  It proves the assertion also at null histories.
For countably many outcomes the sums of positive effects increase
weakly to bounded effects, and the corresponding probabilities
converge by Lemma~\ref{lem:countable91}.  Equivalently retain a remainder
branch and pass to its vanishing trace tail.  Continuity then covers
norm limits of finite testers.  No new quantum system was introduced;
the old system retained by the replacement is precisely $R$, and the
fresh reference remains $F_y$.  Compression to a common support and
extension of the POVM by an arbitrary decision on its complement give
the last assertion.
\end{proof}
A constraint only on the smaller \emph{post-instrument} system is not
sufficient for this resource statement: postponing a compression can
retain a larger $R$.  We do not apply the lemma to bypass such a cut.
When $\sigma_{yh}$ depends on $h$, there need not be one common
$Z_{\theta,y,z}$, so the displayed factorization is unavailable.
First-label dependence $\sigma_y$ causes no such difficulty.

\subsection{An exact criterion for arbitrary finite experiments}
Retain the experiment and the positive payoff matrices
\eqref{eq:gameblocks91}.  Fix a pure first acquisition $|C_0\rangle$
with $C_0:I_1\to R$, $C_0^*C_0=\rho$, and compress $R$ to the range
of $C_0$.  The initial state is held fixed, not replaced by an initially
announced mixture.  Let $1\le\ell\le d_2$ and write
\[
 g_{y,\ell}(a)=
  \max_{b\in\mathcal D_{d_2},\ \operatorname{rank}b\le\ell}\Phi_y(a,b),
 \qquad
 c_{y,\ell}(\rho)=
  \max_{\substack{\rho=\sum_hw_ha_h\,,\ w_h\ge0\\a_h\in\mathcal D_{d_1}}}
                                      \sum_hw_hg_{y,\ell}(a_h).
\]
Here $\Phi_y$ is the final-POVM optimization in
\eqref{eq:leafvalue91}.  The variables $w_h=\tr A_h$ are
\emph{Gram weights}; the physical first-record probability is
$w_h\tr(a_h(E_y^{\theta,1})^{\mathsf T})$.  The two are not
interchangeable.  Both old receivers and their intermediates may be
retained without a bound below $d_1$.  Write
$P_{\mathrm{label},\ell}(\rho)$ and
$P_{\mathrm{msg},\ell}(\rho)$ for the fixed-initial optima when fresh
preparation may use only $y$, or may also use the receiver message.

\begin{theorem}[Uniform message criterion]\label{thm:messagecriterion96}
For every finite experiment, every prescribed pure initial Gram
$\rho$, and every fresh-reference bound $\ell$,
\begin{equation}\label{eq:messagecriterion96}
 \begin{split}
 P_{\mathrm{label},\ell}(\rho)&=\sum_yg_{y,\ell}(\rho),\\
 P_{\mathrm{msg},\ell}(\rho)&=\sum_yc_{y,\ell}(\rho),\\
 \Delta_\ell(\rho)&=
       \sum_y\bigl[c_{y,\ell}(\rho)-g_{y,\ell}(\rho)\bigr]\ge0.
 \end{split}
\end{equation}
Both optima are attained; for the second, at most
$(\operatorname{rank}\rho)^2$ nonzero receiver outcomes after each
first label suffice.  The following statements are equivalent:

\smallskip
\noindent (i) $\Delta_\ell(\rho)=0$ for every prescribed pure first
acquisition $\rho$;

\noindent (ii) every $g_{y,\ell}$ is concave on $\mathcal D_{d_1}$;

\noindent (iii) no protocol with a binary receiver instrument after
one first label, and no receiver instrument after the other labels,
can strictly outperform the optimal label-only value at its own
prescribed initial acquisition.

\smallskip
If messages help somewhere, two outcomes after a single first label
therefore suffice to witness a strict improvement at \emph{some}
initial acquisition.  This last assertion does not fix in advance the
initial acquisition at which the witness occurs.
\end{theorem}
\begin{proof}
Square roots, payoff contractions and the compact final-POVM
optimization make $\Phi_y$ continuous.  The set of rank-at-most-$\ell$
densities is compact, so $g_{y,\ell}$ is continuous and its maximum is
attained.  A mixed fresh preparation with reference dimension $\ell$
may be decomposed into pure vectors on that same space.  With no
receiver message, their mixture is independent of the old record.
Linearity shows that one of these fresh vectors does at least as well
as the mixture for a fixed continuation.  Its Gram has rank at most
$\ell$; conversely every such Gram has a purification with reference
dimension $\ell$.  Lemma~\ref{lem:postponement96} now proves the first
formula without disallowing old quantum information in the final
readout.

With messages, resolve receiver Kraus indices and fresh pure-state
mixtures into the public record.  Every fresh pure coefficient map has
rank at most $\ell$ on its existing reference; refining the receiver
instrument keeps its output space unchanged.  The resulting branch
Grams have the common sum $\rho$ for every $y$, and the best branch
value is $w_hg_{y,\ell}(a_h)$.  Conversely
Lemma~\ref{lem:resetnormal91} implements every finite such decomposition
on the fixed Schmidt support; the fresh factors have rank at most
$\ell$.  The graph of $g_{y,\ell}$ on the density matrices supported
in $\operatorname{supp}\rho$ is compact.  Its convex hull gives an
attaining roof.  If an attaining ensemble has more than
$s^2$ atoms, $s=\operatorname{rank}\rho$, its density matrices are
linearly dependent in a real vector space of dimension $s^2$.
Vary their weights in both directions along this dependence.  An
objective-changing direction would contradict optimality; otherwise
move until a weight vanishes.  Iterating proves the atom bound.
Countable outcomes follow by positive tails.  These arguments prove
the second formula and its attainment.

Each difference in the last line of \eqref{eq:messagecriterion96} is
nonnegative.  If (i) holds, each difference vanishes at every $\rho$;
thus $g_{y,\ell}=c_{y,\ell}$ and (ii) follows.  Conversely concavity
makes each roof equal to its defining function, proving (i).
If (ii) fails, by the definition of concavity there are a label $y_0$,
states $a,b$, and $0<\lambda<1$ such that, for
$\rho=\lambda a+(1-\lambda)b$,
\[
 \lambda g_{y_0,\ell}(a)+(1-\lambda)g_{y_0,\ell}(b)
                                      >g_{y_0,\ell}(\rho).
\]
Realize this two-atom instrument after $y_0$, with its optimizing fresh
continuations.  At every other first label use the optimizing
label-only continuation at $\rho$.  The strict displayed excess is
exactly its excess over the \emph{optimal} label-only value at $\rho$.
This disproves (iii).  Finally (i) implies (iii) by class inclusion.
\end{proof}
The criterion applies to fixed-pair as well as ensemble problems, but
it does not assert that a particular experiment has a positive gap.
It classifies uniform absence of message advantage in the stated
reset interface, not every meaning of bounded coherent memory.
The inclusion diagram at a fixed admissible initial acquisition is
\[
 \begin{array}{ccccc}
 \mathfrak C_\ell(\rho)&\subseteq&\mathfrak R_{d_1,\ell}(\rho)
                                      &\supseteq&\mathfrak L_\ell(\rho),\\
 &&\mathfrak R_{d_1,\ell}(\rho)\subseteq\mathfrak A(\rho).&&
 \end{array}
\]
where $\mathfrak C$ keeps only classical old information,
$\mathfrak L$ allows label feedback and full old quantum storage but
no receiver-to-source message, and $\mathfrak A$ allows coherent
acquisition.  No inclusion between $\mathfrak C$ and $\mathfrak L$
is used.  At a fixed Gram, a message gap is compatible with a globally
optimized no-message strategy using a different first acquisition.

\begin{corollary}[Stability of the message gap]\label{cor:messageperturb96}
Keep the labels, reward table, prior law, and both optimization classes
fixed.  For every latent index $\theta$, suppose the entire slot-$i$
measurement channel changes by at most $\delta_i$ in \emph{unhalved}
diamond norm, $i=1,2$.  The two optima in
\eqref{eq:messagecriterion96} each change by at most
$(\delta_1+\delta_2)/2$, and
\[
       |\widetilde\Delta_\ell(\rho)-\Delta_\ell(\rho)|
                                             \le\delta_1+\delta_2.
\]
In particular a larger original gap remains strict under these
perturbations.  The latent device is still selected once.
\end{corollary}
\begin{proof}
For a fixed protocol, replace the two channel slots in turn.
Diamond-norm control, including every reference, and trace-norm
contractivity bound the change of its final normalized state by
$\delta_1+\delta_2$.  The difference has trace zero, so a reward effect
between zero and the identity changes its expectation by at most
one half of that norm.  Average over the fixed prior and take the
supremum in each unchanged class.  The triangle inequality gives
the gap bound.  The argument compares specified valid reset protocols;
it does not infer physical reset independence from the observed score.
\end{proof}

\section{An exact operational separation of acquisition and receiver memory}
\label{sec:operationalmemory90}
The local profile law does not identify exact operational distances.
We now separate three strategy classes on one finite ensemble of ordered
measurement channels.  The separation holds in every input dimension,
throughout a depolarizing family, with exact optimal success probabilities.
The unknown interface still has only classical output.  The experiment
has a binary hypothesis decision, not a request to identify every member
of the ensemble.

\subsection{A finite family and its second moments}
Let $d\ge2$, let $S$ interchange the two factors of $\C^d\otimes\C^d$,
and put $\Pi_\pm=(I\pm S)/2$.  A weighted family of ordered orthonormal
bases consists of projections $P_y^b$, $1\le y\le d$, and positive
weights $w_b$ summing to one.  The weights are real; rational or algebraic coordinates are not required
by the theorem.  We call this a weighted exact second-moment family,
not necessarily an unweighted projective design.  The required identities are
\begin{equation}\label{eq:basismoments90}
 \sum_b w_b P_y^b\otimes P_z^b=
 \begin{cases}
 (I+S)/(d(d+1)),&y=z,\\
 (dI-S)/(d(d^2-1)),&y\ne z.
 \end{cases}
\end{equation}
\begin{lemma}[Finite simultaneous second moments]\label{lem:finitedesign90}
In every dimension $d\ge2$ there is such a family with at most $d^6+1$
bases.  For $d=2$, the six ordered eigenbases of the three Pauli matrices,
with both orders and equal weights, satisfy \eqref{eq:basismoments90}.
\end{lemma}
\begin{proof}
Average the ordered coordinate basis over Haar measure on $U(d)$.
The second moment of a single column is supported on the symmetric
subspace, invariant under $U\otimes U$, and has trace one.  It is
therefore $2\Pi_+/[d(d+1)]$.  For distinct columns the invariant moment
has the form $aI+bS$; its trace is one and its contraction with $S$ is
zero, since the columns are orthogonal.  These two equations give the
second line of \eqref{eq:basismoments90}.  Equivalently, this line follows
by summing over the other $d-1$ columns and using the first moment $I/d$.
The vector of all $d^2$ Hermitian second moments lies in a real space of
dimension at most $d^6$.  Its Haar average belongs to the convex hull
of the compact set of these vectors.  Carath\'eodory's theorem gives a
convex combination of at most $d^6+1$ vectors; discard zero weights.
For the stated qubit family expand $P=(I\pm\sigma_a)/2$ and use
$S=(I\otimes I+\sum_{a=1}^3\sigma_a\otimes\sigma_a)/2$.
\end{proof}
This is a finite existence argument, not an efficient design construction.
Second-moment designs themselves are established objects \cite{GAE90}.

Fix $0\le t\le1$, and consider the ordered measurements
\begin{equation}\label{eq:gamefamily90}
 C_y=I/d,\qquad E_y^b(t)=(1-t)I/d+tP_y^b.
\end{equation}
Set
\begin{equation}\label{eq:gameprior90}
 L=2(d+1)-t^2,\qquad
 p_C=\frac{d+1-t^2}{L},\qquad p_E=\frac{d+1}{L}.
\end{equation}
With probability $p_C$ the device is $C$; otherwise a basis $b$ is chosen
with law $w_b$ and the device is $E^b(t)$.  This choice is made once.
Both calls use the same chosen, memoryless, classical-output device.
The task is to decide $C$ versus the alternative family.  In particular,
the second moments cannot be replaced by two independent first moments.
Every effect is strictly positive for $t<1$.
Taking the partial trace in \eqref{eq:basismoments90} gives
$\sum_b w_bE_y^b(t)=C_y$.  Thus the two hypotheses have identical
one-call channels, including on entangled inputs, and the one-call
optimum is $p_E$.  The two-call signal concerns a device chosen once;
it is absent if the alternative is independently redrawn at each call.

We compare: (i) complete first-call measurement followed by
measure-and-reprepare classical adaptation, as in
\cite[Definition~23]{OZNPQ89}; (ii) unit-reset protocols of
Definition~\ref{def:reset88}, with profile $(1,1)$ and arbitrary receiver
quantum memory and classical feedback; and (iii) all two-call adaptive
protocols.  Let their optimal Bayes success probabilities be
$P_{\rm ca}$, $P_{\rm reset}$ and $P_{\rm all}$, respectively.
Early stopping is allowed and can be implemented by a discarded fresh call.
The three closed optimization classes are precisely those of
Definition~\ref{def:classes91}.  Lemma~\ref{lem:resetnormal91} gives
the full branch normal form, and Proposition~\ref{prop:causalnormal91}
gives the physical causal normalization independently of this example.

\begin{theorem}[Exact two-call hierarchy]\label{thm:operationalhierarchy90}
The alternative basis index in this experiment is drawn once and
reused at both calls.  Independently redrawing it instead gives the
scalar product channel and no improvement over the larger prior.
For every finite family satisfying \eqref{eq:basismoments90}, every
$d\ge2$, and every $0\le t\le1$, the game \eqref{eq:gamefamily90}--\eqref{eq:gameprior90}
has the exact optimal values
\begin{equation}\label{eq:operationalhierarchy90}
\begin{split}
 P_{\rm ca}&=\frac{d+1}{L},\\
 P_{\rm reset}&=\frac{d+1}{L}+\frac{t^2(d-1)}{2dL},\\
 P_{\rm all}&=\frac{d+1+t^2}{L}.
\end{split}
\end{equation}
For $t>0$ both inequalities are strict.  At $t=0$ all three
values are $1/2$.  The priors sum to one and
$p_E-p_C=t^2/L$, so $p_E$ is the constant-decision baseline.
The middle value is attained
by independent maximally entangled input--reference pairs and a final
joint reference readout, without feedback.  Its upper bound allows all
unit-reset feedback and unrestricted receiver dimension.  The last
value is attained by a parallel antisymmetric two-input state, and its
upper bound allows every adaptive two-call tester.
\end{theorem}
The theorem separates a finite ensemble decision task, not an arbitrary
pair of fixed channels and not a complete memory-dimension hierarchy.
The middle class has hard resources $N=Q=2$; the last construction is
one acquisition group of size two, with square cost four.

\subsection{A swap-filter identity}
For a Hermitian matrix $X$, write $X_+$ for its positive part.
\begin{lemma}[Negative mass of a filtered swap]\label{lem:swapfilter90}
Let $A,B\succeq0$ be $d\times d$ matrices and
$K=(\sqrt A\otimes\sqrt B)S(\sqrt A\otimes\sqrt B)$.  Then
\begin{equation}\label{eq:swapfilter90}
 \tr(-K)_+=\frac12\left[
  \bigl(\tr\sqrt{\sqrt A B\sqrt A}\bigr)^2-\tr(AB)\right]
 \le\frac{d-1}{2d}\tr A\tr B.
\end{equation}
The same negative mass is obtained from $(U\otimes V)S(U\otimes V)^*$
when $U^*U=A$ and $V^*V=B$, even if the receiving spaces are larger.
If $\tr A\tr B>0$, equality in the upper bound holds if and only if
$A=(\tr A)I/d$ and $B=(\tr B)I/d$.
\end{lemma}
\begin{proof}
First suppose $A,B$ are positive definite.  The matrix $K$ is similar
to $(A\otimes B)S$.  Conjugation of the latter by $I\otimes B$ gives
$(AB\otimes I)S$.  The matrix $AB$ is similar to $\sqrt A B\sqrt A$;
let its positive eigenvalues be $\lambda_1,\ldots,\lambda_d$.
Conjugating by the tensor square of an eigenbasis matrix leaves $S$
unchanged.  The resulting matrix has eigenvalues $\lambda_i$ on the
one-dimensional spaces indexed by $(i,i)$, and
$\pm\sqrt{\lambda_i\lambda_j}$ on the two-dimensional spaces indexed
by $(i,j),(j,i)$, $i<j$.  Thus its negative mass is
$\sum_{i<j}\sqrt{\lambda_i\lambda_j}$, which gives the equality.
By Cauchy--Schwarz,
$\sum_i\lambda_i\ge d^{-1}(\sum_i\sqrt{\lambda_i})^2$.
Moreover, the polar decomposition and Hilbert--Schmidt Cauchy--Schwarz
give $\|\sqrt A\sqrt B\|_1\le\sqrt{\tr A\tr B}$.
These inequalities give the bound.  Approximation by positive definite
matrices proves the singular case.  The polar decompositions of $U,V$
identify their filtered swap with $K$ on its support, adding only zero
eigenvalues in the orthogonal complement.
For the equality assertion, put $f=\|\sqrt A\sqrt B\|_1$ and
$a=\tr A$, $b=\tr B$.  Equality in the stated bound, with $ab>0$,
forces equality in both $\tr(AB)\ge f^2/d$ and $f^2\le ab$.
The variational formula for the trace norm and equality in
Hilbert--Schmidt Cauchy--Schwarz imply $B=(b/a)A$: indeed a maximizing
unitary $U$ has $\sqrt B U=c\sqrt A$, and multiplication by its
adjoint gives this proportionality.  Equality in the first inequality
makes all eigenvalues of $\sqrt A B\sqrt A$ equal and positive.
Consequently $A=(a/d)I$ and $B=(b/d)I$.  These scalar matrices plainly
attain the bound.  This also excludes equality at nonzero singular
factors.
\end{proof}
The normalized equality statement and its singular case are expanded
in the proof of Lemma~\ref{lem:centeredswap91}: normalized fidelity one
forces equality of the two density matrices, and eigenvalue
Cauchy--Schwarz then forces all $d$ eigenvalues to be $1/d$.

\subsection{Payoff and complete normalization}
We use the fully transposed, unnormalized Choi representative
\begin{equation}\label{eq:choiconvention90}
 J^\sharp(E)=\left[(\operatorname{Id}\otimes\mathcal M_E)
          (|\Omega\rangle\langle\Omega|)\right]^{\mathsf T}
       =\sum_y E_y\otimes|y\rangle\langle y|,
 \qquad |\Omega\rangle=\sum_i|ii\rangle.
\end{equation}
Its output trace is $I$.  A simultaneous full transpose of the usual
tester representative makes its contraction with $J^\sharp$ the Born
probability.  Factor order is $I_1,Y_1,I_2,Y_2$; expressions grouping
$I_1,I_2$ use the canonical permutation of these factors.
Dephasing the output factors does not change any score.  Write $T_{0,yz}$
for the input block of the tester decision $C$.  The gain over always
guessing the alternative is
\begin{equation}\label{eq:payoff90}
 P-p_E=\sum_{y,z}\tr(T_{0,yz}W_{yz}),\qquad
 W_{yz}=p_C C_y\otimes C_z-p_E\sum_b w_b E_y^b(t)\otimes E_z^b(t).
\end{equation}
The averaged tensor in the second term is $(1-t^2)I/d^2$ plus $t^2$
times \eqref{eq:basismoments90}.  With $k=t^2/(dL)$, substitution gives
\begin{equation}\label{eq:payoffblocks90}
 W_{yy}=-kS,\qquad W_{yz}=-\frac{2k}{d-1}\Pi_-\quad(y\ne z).
\end{equation}
These operators are real in the coordinate basis; a full transpose
therefore does not alter the following congruences.

\begin{proof}[Proof of Theorem~\ref{thm:operationalhierarchy90}]
\emph{Classical adaptation.}
Each block of a measure-and-reprepare tester is a sum of positive
products $A\otimes B$.  Since
$\tr[(A\otimes B)S]=\tr(AB)\ge0$, equal labels contribute no positive
gain in \eqref{eq:payoff90}; different labels cannot contribute positively
either.  Hence $P_{\rm ca}\le p_E$, with equality by the constant decision.
This proof also covers limits of such testers.

\emph{All unit-reset protocols.}
Purify the first acquisition and write its state as
$\sum_i|i\rangle C_0|i\rangle$, with $\tr C_0^*C_0=1$.
After the first device label $y$, combine all intervening receiver
operations into an instrument with recorded outcome $h$.  For an upper
bound, retain its Stinespring environment in the receiver.  The first
conditional reference operator then has the form
$C_{yh}E_y^{\mathsf T}C_{yh}^*$, where
\[
 A_{yh}=C_{yh}^*C_{yh},\qquad
 \sum_h A_{yh}=C_0^*C_0\quad\hbox{for each }y.
\]
Retaining this environment only enlarges receiver processing and does
not import it into the second acquisition.  Conditional reset permits
a purification of that fresh acquisition with a map $D_{yh}$ satisfying
$B_{yh}=D_{yh}^*D_{yh}$ and $\tr B_{yh}=1$.
Its complete system is independent of the old receiver conditional on
$(y,h)$.  The two-reference payoff operator at record $(y,h,z)$ is
therefore
\[
 (C_{yh}\otimes D_{yh})W_{yz}^{\mathsf T}
                         (C_{yh}\otimes D_{yh})^*.
\]
For different labels it is nonpositive.  For equal labels, maximizing
any final effect bounded by the identity gives at most its positive
mass, which by Lemma~\ref{lem:swapfilter90} is
$k(d-1)\tr A_{yh}/(2d)$.  Summing over $y,h$ gives
\[
 P_{\rm reset}-p_E\le k(d-1)/2=\frac{t^2(d-1)}{2dL}.
\]
Public randomization can be included in the recorded history.  A normal
instrument with countably many outcomes is treated by trace-class
sums, or by finite truncations retaining a remainder outcome and then
taking the trace-norm limit.  The bound is independent of receiver
size.  Thus no finite-dimensional restriction on receiver storage has
entered the upper bound.
For $t>0$, the equality case of Lemma~\ref{lem:swapfilter90} also shows
that a nonrandomized protocol attaining this bound in the displayed
purified normal form must satisfy
$A_{yh}=(\tr A_{yh})I/d$ and $B_{yh}=I/d$ on every nonzero branch.
All deficits in the summed upper bound are nonnegative, so equality
forces their branchwise vanishing.  Summing over $h$ then gives
$C_0^*C_0=I/d$.  Maximal mixing is thus a consequence of equality,
not a restriction imposed on the competing protocols.  This statement
concerns the purified normal form, not uniqueness of its dilations
or of the final readout.

For attainment, prepare $|\Omega\rangle/\sqrt d$ independently on each
input and reference.  Declare $C$ precisely when the two device labels
agree and the references lie in $\Pi_-$.  The event probability under
$C$ is $a=(d-1)/(2d^2)$.  Under any $E^b(t)$ it is $(1-t^2)a$:
indeed $\tr E_y^b(t)=1$, $\tr E_y^b(t)^2=t^2+(1-t^2)/d$, and
$\tr[\Pi_-(E\otimes E)]=[(\tr E)^2-\tr E^2]/2$.
The factor $d^{-2}$ comes from the two normalized input--reference
pairs.  The resulting gain is
$a[p_C-p_E(1-t^2)]=t^2(d-1)/(2dL)$.

\emph{All adaptive protocols.}
The deterministic normalization of any two-call tester is
$T_0+T_1=\Xi\otimes I_{Y_2}$, where, after dephasing $Y_1$,
$\Xi$ has input blocks $\Xi_y\succeq0$ satisfying
\[
 \tr_{I_2}\Xi_y=\rho\quad(1\le y\le d),\qquad \tr\rho=1.
\]
Here is a direct realization of these constraints in the present
classical-output interface.  Purify the initial state as
$|\psi\rangle\in I_1\otimes R$, and let $\mathcal A_y$ be the
trace-preserving intervening map from $R$ to $I_2\otimes R'$ after
observing $y$.  Put
\[
 X_y=(\operatorname{Id}_{I_1}\otimes\mathcal A_y)
                 (|\psi\rangle\langle\psi|),\qquad
 \Xi_y=\tr_{R'}X_y,\qquad
 \rho=\tr_R|\psi\rangle\langle\psi|.
\]
Positivity and trace preservation give the displayed marginal equations.
If $F_{yz}$ is the final receiver effect for decision $C$, then
\[
 T_{0,yz}=\tr_{R'}\bigl[(I_{I_1I_2}\otimes\sqrt{F_{yz}})
           X_y(I_{I_1I_2}\otimes\sqrt{F_{yz}})\bigr].
\]
The Choi contraction \eqref{eq:choiconvention90} gives its Born
probability, and $0\preceq F_{yz}\preceq I$ gives
$0\preceq T_{0,yz}\preceq\Xi_y$.  This also proves positivity of the
complement, without an input-dimension factor from a normalized Choi
state.  In particular $\tr\Xi_y=1$.
Since $-kS\preceq kI$ and the other payoff blocks are nonpositive,
\[
 P_{\rm all}-p_E\le k\sum_y\tr\Xi_y=kd=t^2/L.
\]
Prepare any normalized antisymmetric pure state on the two inputs and
declare $C$ when their labels agree.  The event probability is $1/d$
under $C$, and $(1-t^2)/d$ under every $E^b(t)$, because
$(P_y^b\otimes P_y^b)\Pi_-=0$.  The gain is $t^2/L$.
This is a legitimate parallel two-call strategy and attains the
adaptive upper.  The case $t=0$ gives equal values $1/2$ directly.
\end{proof}

\begin{corollary}[Seven qubit devices]\label{cor:qubitgame90}
Take the scalar qubit measurement with prior $2/5$ and each of the six
ordered Pauli measurements with prior $1/10$.  With the device fixed
for both calls and the binary decision specified above, the exact
success probabilities are
\[
 P_{\rm ca}=3/5,\qquad P_{\rm reset}=13/20,\qquad P_{\rm all}=4/5.
\]
\end{corollary}
\begin{proof}
Use Lemma~\ref{lem:finitedesign90} and substitute $d=2,t=1$.
\end{proof}
The antisymmetric readout and comparison mechanism have antecedents in
measurement comparison \cite{ZHS90}; reductions between measurement and
state discrimination also have a substantial history \cite{SZ74}.
Here the assertion is the joint exact optimization over the three
specified classes, including unrestricted reset feedback, on one finite
classical-output ensemble and its full-rank deformation.

\subsection{Quantitative stability of the operational gap}
\begin{corollary}[A stable score separation]\label{cor:robustgame90}
Suppose every device in the finite game is replaced by a classical-output
measurement channel at diamond distance at most $\delta$ from its
specified channel.  Keep the priors and the two-call resource classes
fixed.  Each of their optimal success probabilities changes by at most
$\delta$.  If the displayed reset implementation also has total hard
pathwise preparation defect at most $\varepsilon$ in the sense of
Proposition~\ref{prop:approxreset89}, its score exceeds the optimal
perturbed classical-adaptation score whenever
\begin{equation}\label{eq:robustgap90}
 2\delta+\varepsilon/2<\frac{t^2(d-1)}{2dL}.
\end{equation}
\end{corollary}
\begin{proof}
A two-call hybrid changes each normalized output in unhalved trace norm
by at most $2\delta$.  A fixed decision event changes by at most
$\delta$, and averaging over the unchanged priors preserves this bound.
Taking suprema over a fixed strategy class gives the same two-sided
bound for its optimum.  The preparation certificate changes each
implemented output by at most $\varepsilon$, hence its decision
probability by at most $\varepsilon/2$.  Subtract the classical upper
from the reset lower in \eqref{eq:operationalhierarchy90}.
\end{proof}
These are conditional norm estimates, not physical calibration claims.
In the preceding corollary the same $\delta$ bounds the scalar channel
and each entire basis measurement channel individually.  Priors and the
once-drawn-index law stay fixed; index-dependent channel perturbations
are permitted.  Corollary~\ref{cor:gamestability91} additionally treats
prior and law perturbations, with the unhalved norm factors explicit.
They make the relative scale needed to preserve a small score gap explicit.

\subsection{Hard budget conventions beyond the normalized domain}
\begin{proposition}[Normalization of arbitrary hard thresholds]\label{prop:budgetdomain90}
For positive integer reservations, $N_\pi\le Q_\pi\le N_\pi^2$.
Given real hard bounds $n,q\ge1$, put
\[
 M=\min\{\lfloor n\rfloor,\lfloor q\rfloor\},\qquad
 R=\min\{\lfloor q\rfloor,M^2\}.
\]
The class $\{\pi:N_\pi\le n,\ Q_\pi\le q\}$ equals
$\{\pi:N_\pi\le M,\ Q_\pi\le R\}$, and $M\le R\le M^2$.
Consequently Theorem~\ref{thm:quadraticbudget89} applies with $(M,R)$.
\end{proposition}
\begin{proof}
On each path $\sum n(h)\le\sum n(h)^2\le(\sum n(h))^2$.
Taking hard suprema proves the first assertion.  Bounded integer costs
have integer suprema, so both constraints can be floored.  The first
inequality reduces the call bound to $M$ and the last reduces the
quadratic bound to $R$.  Conversely the reduced bounds imply the original
ones.  Finally $M\le\lfloor q\rfloor$ and $M\le M^2$.
\end{proof}
The normalization $N\le Q$ allows the $N$-singleton lower construction;
it does not say that every admissible individual policy has $Q_\pi\ge N$.
If a hard threshold is below one, only the zero-call class is available.
Fractional-duration reservations and expected-cost constraints are
different models.  None is substituted for a hard integer reservation here.

\section{An exact hierarchy with equal hypothesis priors}
\label{sec:equalprior91}
The strict separation does not require the prior in
\eqref{eq:gameprior90}.  Equal hypothesis priors lead to a different
optimal receiver readout and to another exact hierarchy.  The finite
weighted second-moment family is unchanged.

\subsection{A trace-norm bound}
Let $S$ be the swap on $\C^d\otimes\C^d$ and
$\Pi_\pm=(I\pm S)/2$.
\begin{lemma}[The centered filtered swap]\label{lem:centeredswap91}
For $d\ge2$, $A,B\succeq0$, and $L=\sqrt A\otimes\sqrt B$,
\begin{equation}\label{eq:centeredswap91}
 \|L(I-dS)L\|_1\le\left(d-\frac1d\right)\tr A\tr B.
\end{equation}
For nonzero factors equality holds precisely when $A$ and $B$ are
positive scalar multiples of $I$.  The same norm and bound hold for
rectangular factors having Gram matrices $A,B$.
\end{lemma}
\begin{proof}
Put $K=LSL$ and $f=\|\sqrt A\sqrt B\|_1$.
The spectrum calculation in Lemma~\ref{lem:swapfilter90} gives
$\|K\|_1=f^2$, including singular factors by continuity.
Since $I-dS=2\Pi_--(d-1)S$, the triangle inequality gives
\begin{align*}
 \|L(I-dS)L\|_1
 &\le 2\tr(L\Pi_-L)+(d-1)\|K\|_1\\
 &=\tr A\tr B-\tr(AB)+(d-1)f^2\\
 &\le\tr A\tr B+(d-1-d^{-1})f^2\\
 &\le(d-d^{-1})\tr A\tr B.
\end{align*}
The penultimate step is eigenvalue Cauchy--Schwarz,
$\tr(AB)\ge f^2/d$; the last is $f^2\le\tr A\tr B$.

Here is the equality argument, also supplying a detailed treatment of
singular nonzero factors in Lemma~\ref{lem:swapfilter90}.
Normalize $a=A/\tr A$, $b=B/\tr B$.  The variational formula for the
trace norm selects a unitary $U$ with
$\operatorname{Re}\tr(\sqrt a\sqrt b U)=\|\sqrt a\sqrt b\|_1$.
Both $\sqrt a$ and $\sqrt b U$ have Hilbert--Schmidt norm one.
Equality in Cauchy--Schwarz, with the phase chosen as above, forces
$\sqrt b U=\sqrt a$; multiplication by the adjoint gives $b=a$.
This implication uses no invertibility.  Equality in
$\tr(AB)\ge f^2/d$ then gives $\tr(a^2)=1/d$.  The $d$
nonnegative eigenvalues of $a$, including zeros, sum to one, so equality
in their Cauchy--Schwarz inequality forces all of them to equal $1/d$.
Thus nonzero singular factors cannot attain equality.  Since
$d-1-d^{-1}>0$, equality in \eqref{eq:centeredswap91} forces both
of these equalities.  Scalar factors attain the bound directly from
the eigenvalues $1-d$ on $\Pi_+$ and $1+d$ on $\Pi_-$.
Polar decompositions of rectangular factors preserve the nonzero
spectrum and add only zero eigenvalues, proving the last assertion.
\end{proof}

\subsection{Equal-prior optimization}
\begin{theorem}[Equal-prior operational hierarchy]\label{thm:equalprior91}
Let $(w_b,P_y^b)$ be any weighted exact second-moment family satisfying
\eqref{eq:basismoments90}.  Choose the scalar measurement $C_y=I/d$
with probability $1/2$, or choose the alternative family with probability
$1/2$ and then draw one basis index $b$ with law $w_b$.  In the latter
case the same device $E_y^b(t)=(1-t)I/d+tP_y^b$ is used at both calls.
For $d\ge2$ and $0\le t\le1$, the three classes of
Definition~\ref{def:classes91} have exact optimal success probabilities
\begin{equation}\label{eq:equalprior91}
 \begin{split}
 P_{\rm ca}^{\rm eq}&=\frac12+\frac{t^2(d-1)}{2d(d+1)},\\
 P_{\rm reset}^{\rm eq}&=\frac12+\frac{t^2(d-1)}{2d^2},\\
 P_{\rm all}^{\rm eq}&=\frac12+\frac{t^2}{2d}.
 \end{split}
\end{equation}
Both inequalities are strict for $t>0$, including the full-rank interval
$0<t<1$; all three values equal $1/2$ at $t=0$.
The reset optimum is attained with independent maximally entangled
input--reference pairs and no receiver feedback.  Its upper includes
all unit-reset receiver instruments and arbitrary receiver dimension.
The unrestricted optimum is attained by a normalized antisymmetric
two-input state and bounds every adaptive strategy.
If the alternative basis is instead redrawn independently at each
call, its two-slot averaged channel is the scalar product channel and
all three equal-prior values are $1/2$.
\end{theorem}
\begin{proof}
Put $a=t^2/[2d^2(d+1)]$.  Subtracting the alternative-weighted
second moment from the scalar-weighted product gives
\begin{equation}\label{eq:equalpayoff91}
 W_{yy}=a(I-dS),\qquad
 W_{yz}=-\frac{a}{d-1}(I-dS)\quad(y\ne z).
\end{equation}
Indeed the diagonal moment is
$(1-t^2)I/d^2+t^2(I+S)/[d(d+1)]$, and the off-diagonal
moment is $(1-t^2)I/d^2+t^2(dI-S)/[d(d^2-1)]$.
These substitutions prove \eqref{eq:equalpayoff91}.
Both $S$ and $\Pi_\pm$ are real in the chosen basis; their full
transpose does not alter these payoff blocks.  The gain over the
constant alternative decision is $\sum_{yz}\tr(T_{C,yz}W_{yz})$.

\emph{Classically adaptive upper.}
After a complete first readout, a branch has an input effect
$A\succeq0$, with $\sum_h A_{yh}=\rho_1$ for each $y$ and
$\tr\rho_1=1$.  The second single-call tester has normalization
$B\succeq0$, $\tr B=1$, and decision blocks $0\preceq T_z\preceq B$.
For $\tr A>0$ set $\widehat A=A/\tr A$.
Contracting \eqref{eq:equalpayoff91} with $A$ in the first slot gives
$a\tr A(I-d\widehat A)$ for $z=y$ and its negative divided by
$d-1$ for $z\ne y$.  For a Hermitian $Q$ and $0\preceq T\preceq B$,
$\tr(TQ)\le\tr(BQ_+)$.  Therefore the branch gain is at most
\[
 a\tr A\,\tr[B|I-d\widehat A|]\le a(d-1)\tr A.
\]
The final inequality follows from $0\preceq\widehat A\preceq I$ and
$d\ge2$.  The zero branch contributes zero.
Summing $\sum_{yh}\tr A_{yh}=d$ proves the first upper bound.
It covers arbitrary second-call references and decisions, not only
unassisted classical probes.  Norm closure preserves it.

To attain it, use the same fixed pure input $|u\rangle$ twice and
declare $C$ on unequal labels.  Under $C$ the equal-label probability is
$1/d$; under the averaged alternative it is
\[
 \sum_{by}w_b\langle u,E_y^b(t)u\rangle^2
 =\frac1d+\frac{t^2(d-1)}{d(d+1)}.
\]
This follows directly from the second moments.  The corresponding
Bayes gain is half their difference, namely $ad(d-1)$.

\emph{Unit-reset upper.}
Use Lemma~\ref{lem:resetnormal91} without any restriction on its
rectangular receiver spaces.  On a branch $(y,h)$ put
\[
 X=(C_{yh}\otimes D_{yh})(I-dS)(C_{yh}\otimes D_{yh})^*.
\]
The elementary identity
$\max_{0\preceq F\preceq I}\tr(FX)=\tr(X_+)$ and
\eqref{eq:equalpayoff91} give total branch gain at most
\[
 a\tr(X_+)+a\tr((-X)_+)=a\|X\|_1
 \le a(d-d^{-1})\tr A_{yh},
\]
since $\tr b_{yh}=1$.  There are $d-1$ opposite-sign output blocks,
which cancel their denominator.  Summing over $y,h$ gives
$a(d^2-1)=t^2(d-1)/(2d^2)$.
Countable outcomes and null histories use unnormalized positive tails.

For attainment, prepare two independent normalized maximally entangled
pairs $|\Omega\rangle/\sqrt d$.  On equal device labels declare $C$
when the references lie in $\Pi_-$; on unequal labels declare $C$
when the references lie in $\Pi_+$.  This is one final joint receiver
measurement, with complementary effect $I-F_C$ on each classical
block.  Its event probabilities are
\begin{equation}\label{eq:equalresetevents91}
 q_C=\frac{(d-1)(d+2)}{2d^2},\qquad
 q_b=\frac{(d-1)(d+2-2t^2)}{2d^2}
 \quad\hbox{for every }b.
\end{equation}
To verify them, the receiver block is
$(E_y^b(t))^{\mathsf T}\otimes(E_z^b(t))^{\mathsf T}/d^2$ and
$\tr[\Pi_\pm(A\otimes B)]=[\tr A\tr B\pm\tr(AB)]/2$.
Here $\tr E_y=1$,
$\tr E_y^2=t^2+(1-t^2)/d$, and
$\tr E_yE_z=(1-t^2)/d$ for $y\ne z$.
Counting the $d$ equal-label and $d(d-1)$ unequal-label blocks gives
\eqref{eq:equalresetevents91}.  Half the difference is the reset upper.

\emph{Unrestricted upper.}
By Proposition~\ref{prop:causalnormal91}, $T_{C,yz}\preceq\Xi_y$
and $\tr\Xi_y=1$.  The sum of the positive payoff parts for a fixed
first label is
\[
 (W_{yy})_++\sum_{z\ne y}(W_{yz})_+
 =a(d+1)\Pi_-+a(d-1)\Pi_+\preceq a(d+1)I.
\]
Thus the gain is at most $ad(d+1)=t^2/(2d)$.
Any normalized state supported on the nonzero antisymmetric input
subspace attains it by declaring $C$ on equal labels: their probability
is $1/d$ under $C$ and $(1-t^2)/d$ under every alternative.

The gaps are
\begin{equation}\label{eq:equalgaps91}
 P_{\rm reset}^{\rm eq}-P_{\rm ca}^{\rm eq}
 =\frac{t^2(d-1)}{2d^2(d+1)},\qquad
 P_{\rm all}^{\rm eq}-P_{\rm reset}^{\rm eq}
 =\frac{t^2}{2d^2}.
\end{equation}
Finally $\sum_bw_b\mathcal M_{E^b(t)}=\mathcal M_C$ as a channel
identity.  Two independent draws give its tensor square, hence the
same two-slot statistics even for adaptive testers.  This proves the
last assertion and the endpoint statements.
\end{proof}
The equal priors concern the two hypotheses, not a uniform prior on all
individual devices.  The result removes prior tuning, but still uses
a shared latent device.  It does not assert a gap for every fixed pair.
Together with Theorem~\ref{thm:resetvariational91}, it gives a general
optimization principle and an explicitly solved family, rather than
identifying these two levels of generality.

\subsection{Seven explicit qubit devices}
\label{sec:qubitexplicit91}
For $d=2$ write
\[
 X=\begin{pmatrix}0&1\\1&0\end{pmatrix},\quad
 Y=\begin{pmatrix}0&-i\\i&0\end{pmatrix},\quad
 Z=\begin{pmatrix}1&0\\0&-1\end{pmatrix}.
\]
The devices are $C=(I/2,I/2)$ and, for $A\in\{X,Y,Z\}$ and
$\epsilon\in\{-1,1\}$,
\[
 E^{A,\epsilon}(t)=
 \left(\frac{I+t\epsilon A}{2},\frac{I-t\epsilon A}{2}\right).
\]
The six alternatives have conditional weight $1/6$.
Both outcome orders are included: they cancel the first Pauli moments
and supply the required ordered second moments.  At $t=1$ they are the
six ordered Pauli eigenbasis measurements.  For the prior of
Theorem~\ref{thm:operationalhierarchy90}, $C$ has mass $2/5$ and each
alternative mass $1/10$; for Theorem~\ref{thm:equalprior91}, these masses
are $1/2$ and $1/12$.

For reset acquisition prepare
\[
 |\phi^+\rangle_{I_1R_1}\otimes|\phi^+\rangle_{I_2R_2},\qquad
 |\phi^+\rangle=(|00\rangle+|11\rangle)/\sqrt2.
\]
The inputs are independent.  The two retained references are measured
jointly only after both calls.  Let
$|\psi^-\rangle=(|01\rangle-|10\rangle)/\sqrt2$,
$\Pi_-=|\psi^-\rangle\langle\psi^-|$ and $\Pi_+=I_4-\Pi_-$
on $R_1\otimes R_2$.
The receiver block for any two devices is $E_y^{\mathsf T}\otimes
F_z^{\mathsf T}/4$.  Under $J^\sharp$, the decision tester block is
$F_{C,yz}^{\mathsf T}/4$, where
\[
 F_{C,yz}=\mathbf1_{y=z}\Pi_-
 \quad\hbox{for the prior }2/5,
 \qquad
 F_{C,yz}=\begin{cases}\Pi_-,&y=z,\\\Pi_+,&y\ne z\end{cases}
 \quad\hbox{for equal priors}.
\]
The other block is $(I_4-F_{C,yz})^{\mathsf T}/4\succeq0$.
Their sum is $I_4/4$ for each $(y,z)$, with partial trace
$I_2/2$ over $I_2$ and trace one.  This verifies the causal normalization
in the order $I_1,Y_1,I_2,Y_2$, using the canonical permutation to group
the input copies.  The transpose is required for arbitrary complex
receiver effects, even though the displayed antisymmetric projectors
are real.

At $t=1$, the first reset event has probabilities $1/8$ under $C$ and
zero under every alternative.  Its Bayes score is
$(2/5)(1/8)+(3/5)=13/20$.
The equal-prior reset event has probabilities $1/2$ and $1/4$,
respectively, giving $\frac12(\frac12)+\frac12(\frac34)=5/8$.
For unrestricted acquisition put the normalized state $|\psi^-\rangle$
on $I_1\otimes I_2$, not on the references, and declare $C$ on equal
labels.  The probabilities are $1/2$ and zero, yielding $4/5$ for the
first prior and $3/4$ for equal priors.  For classical adaptation,
the biased-prior optimum is the constant decision $3/5$.
For equal priors, two inputs $|0\rangle$ and the unequal-label event
have probabilities $1/2$ under $C$ and $1/3$ averaged over alternatives,
giving $7/12$.  Consequently the two exact triples on the same devices are
\[
 \frac35<\frac{13}{20}<\frac45,
 \qquad
 \frac7{12}<\frac58<\frac34.
\]
The class upper bounds, not these evaluations alone, prove optimality.

\subsection{An open neighborhood of the equal-prior experiment}
\begin{corollary}[Channel, prior, and ensemble stability]\label{cor:gamestability91}
Fix the family and $t>0$.  Replace the scalar device and each entire
basis measurement channel individually by channels within unhalved
diamond distance $\delta$.  The perturbations may depend on the basis
index and may change all its effects, but must form legal channels on
the same interfaces.  Keep the once-drawn-index rule.  Replace the
scalar prior $1/2$ by $q$ and the conditional basis law $w$ by $\nu$,
and put
\[
 \eta=\tfrac12\sum_b|\nu_b-w_b|,\qquad
 \kappa=\delta+|q-\tfrac12|+\eta.
\]
For each of the three fixed operational classes,
$|\widetilde P_{\rm class}-P_{\rm class}^{\rm eq}|\le\kappa$.
Thus both class gaps remain strict when
$2\kappa<t^2(d-1)/[2d^2(d+1)]$.
If the exhibited reset protocol has, in addition, a uniform
reference-complete diamond certificate of preparation defect whose
hard pathwise sum is at most $\varepsilon$, its implemented score
still exceeds the perturbed classical optimum provided
\[
 2\kappa+\varepsilon/2<\frac{t^2(d-1)}{2d^2(d+1)}.
\]
\end{corollary}
\begin{proof}
For a fixed protocol and device index, the two-call hybrid bounds
output trace distance by $2\delta$.  Normalization makes the change
of any decision probability at most $\delta$, with the factor $1/2$
for unhalved trace distance.  Changing the binary prior changes success
by at most $|q-1/2|$; changing the alternative law changes an average
of $[0,1]$-valued successes by at most $\eta$.  The bound is uniform
in the protocol, so taking suprema preserves it in both directions.
Subtract and use \eqref{eq:equalgaps91}.
The preparation hybrid changes each implemented output by at most
$\varepsilon$, hence its success by at most $\varepsilon/2$.
This last statement concerns a certified implementation of the displayed
reset protocol, not an enlargement of the ideal optimization class.
\end{proof}

\subsection{Quantitative rigidity of near-optimal reset strategies}
The equality condition for the centered swap has a stable form.  It
controls the input Gram matrices in a chosen purified normal form,
rather than specifying unique instruments, dilations, or receiver
readouts.  For $d\ge2$ set
\[
 c_d=d-1-d^{-1}>0,\qquad
 K_d=d\left(1+\frac{9}{2c_d}\right).
\]
\begin{lemma}[Stable centered-swap equality]\label{lem:stableswap91}
For density matrices $a,b$ on $\C^d$, let
\[
 \Delta(a,b)=d-d^{-1}
  -\| (\sqrt a\otimes\sqrt b)(I-dS)
                         (\sqrt a\otimes\sqrt b)\|_1.
\]
Then $\Delta(a,b)\ge0$ and
\begin{equation}\label{eq:stableswap91}
 \|a-I/d\|_1^2\le K_d\Delta(a,b),\qquad
 \|b-I/d\|_1^2\le K_d\Delta(a,b).
\end{equation}
No invertibility is required.
\end{lemma}
\begin{proof}
Put $f=\|\sqrt a\sqrt b\|_1$, $\eta=1-f^2$, and
$e=\tr(ab)-f^2/d$.  The proof of Lemma~\ref{lem:centeredswap91}
gives $\eta,e\ge0$ and
\begin{equation}\label{eq:swapdefectdecomposition91}
 \Delta(a,b)\ge c_d\eta+e.
\end{equation}
Choose a unitary $V$ extending the right polar factor so that
$Z=\sqrt a\sqrt b V=\sqrt{\sqrt a b\sqrt a}\succeq0$.
This is possible also when the matrices are singular.  The
Hilbert--Schmidt norm, denoted by $\|\cdot\|_2$, satisfies
\[
 \|\sqrt b V-\sqrt a\|_2^2=2(1-f),\quad
 \|Z-a\|_2\le\sqrt{2\eta},\quad
 \|Z-fI/d\|_2^2=e.
\]
Here $\|\sqrt a\|_{\rm op}\le1$ and $1-f\le\eta$.
Since also $1-f\le\sqrt\eta$, the triangle inequality yields
\[
 \|a-I/d\|_2\le(\sqrt2+d^{-1/2})\sqrt\eta+\sqrt e.
\]
Weighted Cauchy--Schwarz bounds its square by
$[1+(\sqrt2+d^{-1/2})^2/c_d](c_d\eta+e)$.
For $d\ge2$, $(\sqrt2+d^{-1/2})^2\le9/2$; and
$\|T\|_1^2\le d\|T\|_2^2$ for a $d\times d$ matrix $T$.
These facts and \eqref{eq:swapdefectdecomposition91} prove the first
estimate.  Interchanging $a,b$ proves the second, since swap
conjugation leaves the centered-swap norm unchanged.
\end{proof}

\begin{corollary}[Near-optimal acquisition marginals]\label{cor:resetstability91}
Fix the equal-prior experiment of Theorem~\ref{thm:equalprior91}
with $t>0$.  Suppose a unit-reset strategy achieves reward
$v\ge P_{\rm reset}^{\rm eq}-\varepsilon$.
Choose its purified normal form from Lemma~\ref{lem:resetnormal91}.
On nonzero branches put
\[
 a_{yh}=A_{yh}/\tr A_{yh},\qquad
 \mu_{yh}=\tr A_{yh}/d.
\]
Then $\sum_{yh}\mu_{yh}=1$ and
\begin{equation}\label{eq:resetstability91}
 \sum_{y,h}\mu_{yh}
 \left(\|a_{yh}-I/d\|_1^2+\|b_{yh}-I/d\|_1^2\right)
 \le \frac{4K_d d(d+1)}{t^2}\,\varepsilon.
\end{equation}
For a public mixture the same conclusion holds after adjoining the
seed and multiplying these weights by its law.  Finite and countable
recorded instruments are both allowed.
\end{corollary}
\begin{proof}
Write $\alpha=t^2/[2d^2(d+1)]$ for the coefficient in
\eqref{eq:equalpayoff91}.  The complete reset upper, before applying
the centered-swap bound, gives
\[
 v\le\tfrac12+\alpha\sum_{yh}\tr A_{yh}
          [d-d^{-1}-\Delta(a_{yh},b_{yh})].
\]
Purification does not obstruct recovering the original decision rule,
so this inequality bounds its achieved reward as well.  The common
Gram identity gives $\sum_{yh}\tr A_{yh}=d$.  Subtract from the
attained optimum to obtain
$\sum_{yh}\mu_{yh}\Delta(a_{yh},b_{yh})
\le\varepsilon/(\alpha d)$.
Apply Lemma~\ref{lem:stableswap91} to both factors and substitute
$\alpha$.  Positive trace tails justify countable sums.  For a public
mixture, apply the same upper to each component and average before
subtracting its reward from the common optimum.
\end{proof}
In particular, the total $\mu$-weight on which either normalized Gram
matrix is at trace distance at least $u>0$ from $I/d$ is at most
$4K_d d(d+1)\varepsilon/(t^2u^2)$.  These are Gram weights, not
probabilities under a selected hypothesis.  At zero loss the result
recovers branchwise isotropy.  The estimate is nontrivial at the
relative scale $\varepsilon/t^2$; it does not turn fixed calibration
error into vanishing-signal robustness.  It asserts neither uniqueness
of a dilation nor a device-independent certification of the reset cut.

\section{Rank-resolved receiver resources}
\label{sec:rankmemory92}
The distinction between classical and quantum receivers extends to an
attained hierarchy at each pair of retained-register dimensions.  This
section keeps the reset cut fixed and resolves its two sides separately.
It supplies a general variational characterization and an exactly solved
spectrum, rather than identifying all notions of bounded quantum memory.

\subsection{Two specified recording cuts}
Use the finite experiment of Definition~\ref{def:classes91}.
Fix integers $1\le k\le d_1$ and $1\le \ell\le d_2$.
\begin{definition}[Retained-register dimensions]\label{def:rankclass92}
The class $\mathfrak R_{k,\ell}$ consists of unit-reset strategies for
which, conditional on every recorded first history, the entire retained
old quantum system has dimension at most $k$ and the entire fresh
reference accompanying the second input has dimension at most $\ell$.
The old system and the complete fresh input--reference system are in a
tensor product at that cut.  All other quantum registers are discarded
before the second call; classical records may be retained without a
cardinality bound.  The initial first-call reference and temporary
workspaces before this cut have no dimension restriction.  The final
measurement on the two retained receivers is arbitrary.
Mixed preparations, recorded normal instruments, public randomization,
null histories, discarded terminal calls and finite-dimensional tester
closure have the conventions of Definition~\ref{def:classes91}.
Denote the optimal reward by $P_{k,\ell}$.
\end{definition}
In particular, this is not a bound on the largest workspace throughout
an implementation.  Enlarging the old output by keeping an instrument
environment would violate the stated resource.  The refinement in the
next lemma records that environment's Kraus index \emph{classically}
instead, without retaining it coherently.

\begin{lemma}[Dimension-preserving refinement]\label{lem:ranknormal92}
For the optimization over $\mathfrak R_{k,\ell}$, the normal form
\eqref{eq:resetgrams91}--\eqref{eq:resetblock91} is exact with
\begin{equation}\label{eq:rankconstraints92}
 \operatorname{rank}A_{yh}\le k,\qquad
 \operatorname{rank}b_{yh}\le\ell.
\end{equation}
Every implemented competitor can be refined into such branches while
preserving its score if the added classical indices are ignored in
its decisions.  Conversely every finite Gram collection satisfying
these ranks and the common-barycenter constraints is realizable in
$\mathfrak R_{k,\ell}$.
\end{lemma}
\begin{proof}
Purify the first input, retaining a reference whose relevant Schmidt
support has dimension at most $d_1$.  A completely positive branch of
the subsequent map from that support to $\C^k$ has a finite Kraus
representation $X\mapsto\sum_qV_{yhq}XV_{yhq}^*$, with
$V_{yhq}:\C^{r_0}\to\C^k$.  Refine $h$ to $(h,q)$, record $q$
classically, and use the original fresh preparation and final effects.
Summing over $q$ recovers the old channel and its decision statistics.
Thus $C_{yhq}=V_{yhq}C_0$ has rank at most $k$, and completeness
still gives $\sum_{h,q}C_{yhq}^*C_{yhq}=\rho$ for every $y$.
An original mixed first preparation is recovered by discarding its
purification before this old-output cut; this discarding is part of
that same completely positive branch, not an additional retained
quantum register.

A mixed fresh state on $I_2\otimes\C^\ell$ has a finite spectral
mixture of normalized pure states in that same space.  Sample and
record its mixture index, use the corresponding vector, and ignore
that index in the original decision.  Its coefficient map
$D:I_2\to\C^\ell$ has Gram rank at most $\ell$.  Multiplying the
old coefficient map by the square root of this mixing probability
preserves the common Gram sum.  This is a classical refinement, not a
purification into a larger fresh reference.  It is independent of the
old quantum output conditional on the refined record.  Countable
outcomes are handled by the trace tails of Lemma~\ref{lem:countable91}.
Public mixtures and tester limits preserve the resulting value bound.

Conversely choose $C_0^*C_0=\rho$ on its Schmidt support.  Factor each
$A_{yh}$ as $C_{yh}^*C_{yh}$ with
$C_{yh}:I_1\to\C^k$, and each $b_{yh}$ as $D_{yh}^*D_{yh}$ with
$D_{yh}:I_2\to\C^\ell$.  The maps $V_{yh}=C_{yh}C_0^+$ satisfy
$\sum_hV_{yh}^*V_{yh}=I$ on that support.  They give a recorded
trace-preserving instrument.  The normalized fresh vector $|D_{yh}\rangle$
is prepared independently, and polar isometries transfer the prescribed
leaf POVMs to $\C^k\otimes\C^\ell$.  Zero Gram branches are simply
zero maps.  The original unnormalized formulas define all null histories.
\end{proof}

\subsection{An attained hierarchy and its dual}
Write $\mathcal D_d^{\le q}$ for the compact set of density matrices of
rank at most $q$.  In the notation of \eqref{eq:leafvalue91}, put
\begin{align}
 g_y^\ell(a)&=\max_{b\in\mathcal D_{d_2}^{\le\ell}}\Phi_y(a,b),
 \label{eq:rankleaf92}\\
 c_y^{k,\ell}(\rho)&=\max\left\{\sum_h\lambda_hg_y^\ell(a_h):
 a_h\in\mathcal D_{d_1}^{\le k},\quad
 \lambda_h\ge0,\quad\sum_h\lambda_ha_h=\rho\right\}.
 \label{eq:rankroof92}
\end{align}
Weights sum to one by taking traces.  All density matrices have such a
representation, since rank-one atoms are allowed.

\begin{theorem}[Rank-resolved variational principle]\label{thm:rankroof92}
For every finite two-call experiment,
\begin{equation}\label{eq:rankprimaldual92}
\begin{split}
 P_{k,\ell}&=\max_{\rho\in\mathcal D_{d_1}}
                       \sum_yc_y^{k,\ell}(\rho)\\
 &=\min_{(H_y)}\lambda_{\max}\left(\sum_yH_y\right),\qquad
 \tr(H_ya)\ge g_y^\ell(a)\quad(a\in\mathcal D_{d_1}^{\le k}),
\end{split}
\end{equation}
where the $H_y$ are Hermitian.  Both extrema are attained.  At an
optimal barycenter $\rho$, at most $(\operatorname{rank}\rho)^2$
nonzero receiver outcomes after each first label suffice; a final
quantum receiver of dimension $k\ell$ suffices.
The values increase monotonically in either dimension, with
\[
 P_{1,d_2}=P_{\rm ca},\qquad P_{d_1,d_2}=P_{\rm reset}.
\]
The contact, spectral and quantitative defect criteria of
Theorem~\ref{thm:resetdual91} and Corollary~\ref{cor:contactdefect91}
hold with $g_y^\ell$ and the rank-restricted atoms.
\end{theorem}
\begin{proof}
The function $\Phi_y$ is continuous and homogeneous in its first
positive argument.  The rank-bounded sets are compact, so
$g_y^\ell$ is continuous and its graph on
$\mathcal D_{d_1}^{\le k}$ has compact convex hull.
The largest height over each $\rho$ is an attained, upper
semicontinuous concave envelope.  Lemma~\ref{lem:ranknormal92} gives
both operational directions.  The atom-elimination proof of
Theorem~\ref{thm:resetvariational91} uses only linear dependence of
atoms on the support of $\rho$, and does not change any atom.  It
therefore preserves the rank restriction and gives the same branch
bound.  Compactness supplies all fresh states and final POVMs.

For clarity, dual attainment does not require the rank-bounded atom
set to be convex.  The homogeneous hypograph of its concave envelope
is a closed convex cone, finite on every positive matrix and
nonnegative there.  At the common matrix $I/d_1$ with every height
$-1$, all these cones are strictly feasible.  The conic argument of
Theorem~\ref{thm:resetdual91} applies without alteration: its dual is
precisely majorization of the atom graph in \eqref{eq:rankprimaldual92}.
Bounded reward gives a finite value.  Slater duality yields no gap and
an attained Hermitian majorant.  Subtracting the primal value from the
dual upper gives the same three nonnegative defects, proving the
contact and quantitative assertions.  Inclusion of the rank-bounded
atom sets proves monotonicity.  The two endpoint identifications are
Theorems~\ref{thm:classicalroof91} and~\ref{thm:resetvariational91}.
\end{proof}
The dual inequalities are universal on their rank-bounded domains;
sampling is not a certificate.  In particular, strict improvement
between any two nested dimension pairs is equivalent to a feasible
finite strategy for the larger pair exceeding an attained dual upper
for the smaller pair.  This criterion covers arbitrary fixed pairs,
without asserting that every such pair exhibits a gap.

\subsection{Rank-sensitive swap bounds}
The rank in the next lemma is an upper bound on each of the two Gram
ranks; the active overlap may have a smaller rank.
\begin{lemma}[Swap bounds at finite rank]\label{lem:rankswap92}
Let $1\le r\le d$, $d\ge2$, and let $A,B\succeq0$ have ranks at most
$r$.  With $L=\sqrt A\otimes\sqrt B$,
\begin{align}
 \tr[-LSL]_+&\le\frac{r-1}{2r}\tr A\tr B,
 \label{eq:rankswapnegative92}\\
 \|L(I-dS)L\|_1&\le\left(d-\frac1r\right)\tr A\tr B.
 \label{eq:rankswapcenter92}
\end{align}
The bounds also hold if only the smaller of the two ranks is at most
$r$, and for rectangular factors with the same Grams.  Both bounds
are sharp: take $A=aP/r$, $B=bP/r$, where $P$ is any rank-$r$
projection and $a,b\ge0$.  For nonzero factors, equality in the first
bound with $r>1$ requires these forms with a common $P$; equality in
the second requires them when $(d,r)\ne(2,1)$.
\end{lemma}
\begin{proof}
The nonzero eigenvalues of $\sqrt A B\sqrt A$ number at most $r$.
Consequently, for $f=\|\sqrt A\sqrt B\|_1$,
\[
 \tr(AB)\ge f^2/r,\qquad f^2\le\tr A\tr B.
\]
Insert the first inequality into Lemma~\ref{lem:swapfilter90} to get
\eqref{eq:rankswapnegative92}.  The triangle step of
Lemma~\ref{lem:centeredswap91} gives
\[
 \|L(I-dS)L\|_1
 \le\tr A\tr B-\tr(AB)+(d-1)f^2
 \le\tr A\tr B+(d-1-r^{-1})f^2.
\]
The coefficient is nonnegative, proving \eqref{eq:rankswapcenter92}.
These arguments only use the smaller rank, not a common support
assumption.  The spectrum identities hold for singular factors by the
proved continuity argument, and polar isometries handle rectangular
maps.

For equality, when the coefficient of $f^2$ is positive, normalized
fidelity one gives $A/\tr A=B/\tr B$ without an invertibility
assumption, as in Lemma~\ref{lem:centeredswap91}.  Equality in the
rank-$r$ Cauchy--Schwarz inequality then forces exactly $r$ equal
nonzero eigenvalues, proving the common-projection forms.
For such forms the centered spectrum is $(1-d)ab/r^2$ with
multiplicity $r(r+1)/2$ and $(1+d)ab/r^2$ with multiplicity $r(r-1)/2$;
the negative swap has mass $ab(r-1)/(2r)$.  This proves sharpness.
The excluded centered case has zero coefficient and different equality
cases: for $d=2,r=1$, orthogonal pure factors also attain the bound.
For the first inequality at $r=1$, the negative mass is always zero.
\end{proof}

\subsection{The complete two-cut spectrum of the basis experiment}
The following formula resolves every pair of retained-register
thresholds on the same once-selected finite basis family.  No
additional design is introduced at a new threshold.
\begin{theorem}[Exact dimension spectra]\label{thm:rankspectrum92}
Fix $d\ge2$, $0\le t\le1$, and a weighted exact second-moment family
satisfying \eqref{eq:basismoments90}.  Under the once-drawn device rule,
let $P_{k,\ell}^{\rm b}$ use the priors of
Theorem~\ref{thm:operationalhierarchy90}, and let
$P_{k,\ell}^{\rm eq}$ use the equal hypothesis priors of
Theorem~\ref{thm:equalprior91}.  For $1\le k,\ell\le d$, set
$r=\min\{k,\ell\}$ and $L_0=2(d+1)-t^2$.  Then
\begin{align}
 P_{k,\ell}^{\rm b}
 &=\frac{d+1}{L_0}+\frac{t^2(r-1)}{2rL_0},
 \label{eq:rankbiased92}\\
 P_{k,\ell}^{\rm eq}
 &=\frac12+\frac{t^2(dr-1)}{2dr(d+1)}.
 \label{eq:rankequal92}
\end{align}
These are optima over the complete classes
$\mathfrak R_{k,\ell}$.  For $t>0$ they strictly increase at every
increase of the smaller threshold, and attain the unrestricted reset
values precisely at $k=\ell=d$.  The two-cut resource dependence is
therefore exactly through $\min\{k,\ell\}$ for these experiments.
\end{theorem}
\begin{proof}
Use the refined normal form of Lemma~\ref{lem:ranknormal92}; there is
no retained dilation environment outside the counted receivers.
Every branch has $\tr b_{yh}=1$ and smaller Gram rank at most $r$,
while $\sum_{yh}\tr A_{yh}=d$.
For the biased priors the equal-label payoff is
$-t^2 S/(dL_0)$ and all unequal-label payoffs are nonpositive.
Equation~\eqref{eq:rankswapnegative92} bounds the gain by
$t^2(r-1)/(2rL_0)$.
For equal priors write $\alpha=t^2/[2d^2(d+1)]$.
The $d-1$ opposite-sign payoff blocks cancel their denominator, so the
branch gain is at most
$\alpha\|(C_{yh}\otimes D_{yh})(I-dS)(C_{yh}\otimes D_{yh})^*\|_1$.
Equation~\eqref{eq:rankswapcenter92} and the Gram sum give
$\alpha d(d-1/r)$, which is \eqref{eq:rankequal92} minus $1/2$.
Classical mixtures and trace-norm closure preserve both upper bounds.

Here is a single finite construction attaining them.  Index the input
basis cyclically modulo $d$, and let $J_h:\C^r\to\C^d$ include the
$r$ consecutive coordinates starting at $h$, $0\le h<d$.
Put $P_h=J_hJ_h^*$; then $\sum_hP_h=rI$.
Choose the first coefficient map $C_0=I/\sqrt d$ and, after every
first device label, the receiver instrument and fresh coefficient maps
\begin{equation}\label{eq:cyclicrankprotocol92}
 V_h=J_h^*/\sqrt r,\qquad
 C_{yh}=J_h^*/\sqrt{dr},\qquad D_{yh}=J_h^*/\sqrt r.
\end{equation}
Thus $\sum_hV_h^*V_h=I$,
$A_{yh}=P_h/(dr)$, $b_{yh}=P_h/r$ and $\sum_hA_{yh}=I/d$.
The fresh vector has norm one and is independent of the old receiver.
Both retained quantum registers have dimension $r$; $h$ is classical.
The initial reference has dimension $d$, as explicitly allowed before
the recording cut in Definition~\ref{def:rankclass92}.

On these receivers the filtered swap is $S_r/(dr^2)$ and the centered
one is $(I_{r^2}-dS_r)/(dr^2)$.  For the biased prior, declare $C$
only on equal labels and the antisymmetric receiver projection.
For equal priors, declare $C$ on the antisymmetric projection on
equal labels and on the symmetric projection on unequal labels.
Complementary effects complete a POVM.  Their spectral signs attain
every branch bound above.  At $r=1$ the antisymmetric projection is
zero; the second readout reduces to the unequal-label rule.  At $r=d$
the instrument's redundant classical randomization may be omitted.
All constructions remain normalized at $t=0$ and on projective null
histories at $t=1$.
\end{proof}

\begin{corollary}[Adjacent gains and dimension certificates]
\label{cor:rankwitness92}
For $1\le r<d$, abbreviate the values in
Theorem~\ref{thm:rankspectrum92} by $P_r^{\rm b}$ and $P_r^{\rm eq}$.
Their adjacent gains are
\[
 P_{r+1}^{\rm b}-P_r^{\rm b}=\frac{t^2}{2L_0r(r+1)},\qquad
 P_{r+1}^{\rm eq}-P_r^{\rm eq}
       =\frac{t^2}{2d(d+1)r(r+1)}.
\]
For the equal-prior perturbations of Corollary~\ref{cor:gamestability91},
with its score error $\kappa$, every perturbed class value moves by at
most $\kappa$.  In particular an achieved ideal-reset score strictly
above $P_r^{\rm eq}+\kappa$ requires both counted receiver dimensions
to exceed $r$.  The adjacent gap survives if $2\kappa$ is smaller
than its displayed value.  A certified additional preparation defect
$\varepsilon$ changes the implemented score by at most $\varepsilon/2$
and may be added to these error budgets.
\end{corollary}
\begin{proof}
Subtract the consecutive exact formulas.  The hybrid and probability
law bounds are uniform in a fixed protocol, independently of its
receiver dimensions, so they hold for every rank-restricted supremum.
If either dimension were at most $r$, its nominal value would be at
most $P_r^{\rm eq}$ by Theorem~\ref{thm:rankspectrum92}.  This proves
the asserted certificate and the perturbation statements.
\end{proof}
These dimension certificates assume the specified reset architecture
and calibration bounds.  They are not device-independent conclusions
from an uncalibrated score.  The entire two-cut hierarchy still lies
inside the reset class; the unrestricted coherent-acquisition optimum
of Theorem~\ref{thm:equalprior91} is a different resource.

\subsection{Relation to constrained-memory discrimination}
Memory-constrained tester optimization and its relation to constrained
separability are established in \cite{OZNPQ89}, including the classical
adaptive boundary.  Theorem~\ref{thm:rankroof92} is a rank-envelope
characterization for the specified classical-output, fresh-preparation
cut, and Theorem~\ref{thm:rankspectrum92} solves its two-coordinate
values on one fixed finite family.  It does not replace the general
channel-memory formulations of that work.  The autonomous multi-time
hierarchy of \cite{ZB90} carries coherent information between successive
acquisitions and compiles arbitrary finite-time testers with a counter.
Here increasing the two receiver dimensions never removes the reset
cut.  Consequently saturation means the full reset optimum, not the
unrestricted strategy norm.  These different saturation statements
are compatible and cannot be interchanged.

\section{Initial spectra and rigidity of rank-optimal receivers}
\label{sec:spectralrigidity93}
The dimension spectrum determines an optimal value.  We now determine
which initial spectra can attain it, and how a small loss in value
restricts the acquisition marginals.  Throughout this section the basis
index is drawn once and retained for both calls, the hypothesis priors
are equal, and $0<t\le1$.  The family satisfies
\eqref{eq:basismoments90}.  Put
\[
 r=\min\{k,\ell\},\qquad
 p_r=\frac12+\frac{t^2(dr-1)}{2dr(d+1)},\qquad
 \alpha=\frac{t^2}{2d^2(d+1)}.
\]
All ranks and norms refer to the fixed input copies.  Trace norms are
unhalved.  An initial Gram matrix $\rho$ means a pure first acquisition
$|C_0\rangle$ with $C_0^*C_0=\rho$, in the convention of
Lemma~\ref{lem:resetnormal91}; its input marginal is $\rho^{\mathsf T}$.
The receiver cut is exactly Definition~\ref{def:rankclass92}.

\subsection{The complete set of optimal initial spectra}
\begin{theorem}[Optimal initial Gram matrices]\label{thm:initialspectra93}
A fixed initial Gram matrix $\rho\in\mathcal D_d$ admits a strategy
in $\mathfrak R_{k,\ell}$ attaining $p_r$ if and only if
\begin{equation}\label{eq:spectralcap93}
                  \rho\preceq I/r.
\end{equation}
Every such $\rho$ has a decomposition
\begin{equation}\label{eq:projectionframe93}
 \rho=\sum_{h=1}^{s}w_hP_h/r,\qquad s\le d,\quad
 w_h>0,\quad\sum_hw_h=1,
\end{equation}
where the $P_h$ are rank-$r$ projections commuting with $\rho$.
This decomposition gives an attaining strategy with the same receiver
instrument after every first device label and with both retained
quantum registers of dimension $r$.

Except when $(d,r)=(2,1)$, every attaining strategy in the
dimension-preserving refined normal form has, on each nonzero branch,
\begin{equation}\label{eq:optimalgrams93}
       A_{yh}=w_{yh}P_{yh}/r,\qquad b_{yh}=P_{yh}/r,
       \qquad \operatorname{rank}P_{yh}=r.
\end{equation}
Together with the common Gram sums and maximizing leaf decisions,
these conditions are also sufficient.  They specify Gram data, not
unique dilations or unique final measurements.
\end{theorem}
\begin{proof}
For a normalized branch write $a=A/\tr A$ and $b=b_{yh}$.
The equal-prior upper in Theorem~\ref{thm:rankspectrum92}, before
applying its final inequality, is
\begin{equation}\label{eq:unsaturatedrankupper93}
 v\le\frac12+\alpha\sum_{y,h}w_{yh}
       \| (\sqrt{a_{yh}}\otimes\sqrt{b_{yh}})(I-dS)
                    (\sqrt{a_{yh}}\otimes\sqrt{b_{yh}})\|_1.
\end{equation}
The smaller Gram rank is at most $r$, and $\sum_{y,h}w_{yh}=d$.
All branch deficits from the bound $d-1/r$ are nonnegative.
Attainment therefore forces equality on every branch of positive
weight.  When $(d,r)\ne(2,1)$, Lemma~\ref{lem:rankswap92} gives
\eqref{eq:optimalgrams93}.  For each fixed $y$, summing yields
$\rho=\sum_hw_{yh}P_{yh}/r\preceq I/r$, since
$\sum_hw_{yh}=1$.  In the exceptional case $r=1$, the claimed spectral
condition holds for every density matrix anyway.  The fixed-$\rho$
compact-envelope argument of Theorem~\ref{thm:rankroof92} gives
attainment of its constrained value, so the same conclusion applies
when an optimum was initially defined through tester closure.

We give a finite construction for the converse.  Diagonalize $\rho$
with eigenvalues $\lambda_i$ and set $x_i=r\lambda_i$ and
$s_i=\sum_{j\le i}x_j$, $s_0=0$.  Then $0\le x_i\le1$ and
$s_d=r$.  For almost every $u\in[0,1)$, mark the $r$ points
$u,u+1,\ldots,u+r-1$ in $[0,r)$.  Select coordinate $i$ if its
half-open interval $[s_{i-1},s_i)$ contains a marked point.
Each interval has length at most one, so exactly $r$ distinct
coordinates are selected.  The probability of selecting coordinate
$i$ is $x_i$.  The selected subset is constant between successive
fractional parts of the $s_i$.  There are at most $d$ such intervals
of positive length, because $s_0$ and $s_d$ have the same fractional
part.  Their lengths are the weights $w_h$, and their selected
coordinate projections are $P_h$.  Taking diagonal entries proves
\eqref{eq:projectionframe93}.  Endpoints have zero measure and may
be assigned either adjacent choice.  Rational eigenvalues give rational
weights by this construction; no rational eigenbasis is presumed.

Choose isometries $J_h:\C^r\to\C^d$ with $J_hJ_h^*=P_h$ and put
\[
 C_{yh}=\sqrt{w_h/r}\,J_h^*,\qquad
 D_{yh}=J_h^*/\sqrt r,\qquad V_{yh}=C_{yh}C_0^+.
\]
On the initial Schmidt support, the common Gram sum gives
$\sum_hV_{yh}^*V_{yh}=I$ and $V_{yh}C_0=C_{yh}$.
Thus $V_{yh}$ is a recorded instrument, and $|D_{yh}\rangle$ is an
independently prepared normalized fresh vector.  The same maps may be
used for every $y$.  The receiver centered swap is
$w_h(I_{r^2}-dS_r)/r^2$.  Declare the scalar hypothesis on the
antisymmetric receiver projection for equal labels and on the symmetric
projection for unequal labels.  Complementary effects complete the
measurement.  These decisions attain each positive-part bound, proving
sufficiency, including $r=1$.  Applying the same construction without
a prescribed commuting eigenbasis proves sufficiency of
\eqref{eq:optimalgrams93} with the common Gram sums.  Countable
normal forms are covered by the positive trace tails of
Lemma~\ref{lem:countable91}.
\end{proof}
In particular the initial state need not be maximally mixed when
$r<d$.  Its largest Schmidt weight must be at most $1/r$.
When $r=d$, the condition reduces to $\rho=I/d$.
The finite interval construction concerns a known initial spectrum,
not construction of the second-moment family itself.

\subsection{Counting the first acquisition as well}
Let $\mathfrak R_{q;k,\ell}$ impose the additional condition that the
entire reference accompanying the first input has dimension at most
$q$, conditional on the public preparation record.  Mixed states on
$I_1\otimes\C^q$ are allowed.  The old receiver after its recorded
instrument has dimension at most $k$ and the independent second
reference at most $\ell$, as before.  Workspaces discarded before a
specified cut do not enlarge a counted register.  Classical records
and the final joint receiver measurement have the same conventions as
Definition~\ref{def:rankclass92}.
\begin{corollary}[Three specified dimension cuts]\label{cor:threecuts93}
For $1\le q,k,\ell\le d$, put $s=\min\{q,k,\ell\}$.
On the same finite family the complete classes
$\mathfrak R_{q;k,\ell}$ have exact optimal values
\begin{equation}\label{eq:threecuts93}
 \begin{split}
 P_{q;k,\ell}^{\rm eq}
      &=\frac12+\frac{t^2(ds-1)}{2ds(d+1)},\\
 P_{q;k,\ell}^{\rm b}
      &=\frac{d+1}{2(d+1)-t^2}
          +\frac{t^2(s-1)}{2s[2(d+1)-t^2]}.
 \end{split}
\end{equation}
The second line uses exactly the prior of
Theorem~\ref{thm:operationalhierarchy90}.  Both optima are attained by
two independently prepared maximally entangled vectors on any fixed
rank-$s$ input subspace, with no intervening receiver instrument and
no feedback.  No first-call reference of dimension $d$ is required
unless $s=d$.
\end{corollary}
\begin{proof}
Resolve a mixed first preparation into a finite pure-state mixture in
$I_1\otimes\C^q$ and record its mixture index classically.  This
recovers its old score if the index is ignored, without adding a quantum
purification.  For each component $\operatorname{rank}C_0\le q$.
The Kraus refinement of Lemma~\ref{lem:ranknormal92} then gives
$\operatorname{rank}A_{yh}\le\min\{q,k\}$, while
$\operatorname{rank}b_{yh}\le\ell$.  The two rank-sensitive swap
bounds apply with $s$.  The trace sums, constant-decision baselines
and positive-part optimizations in Theorem~\ref{thm:rankspectrum92}
give both upper bounds.  Averaging over initial mixtures and public
seeds, and then taking norm limits, preserves these bounds.

For attainment fix an isometry $J:\C^s\to\C^d$ and use
$C_0=D=J^*/\sqrt s$.  Keep the first receiver unchanged and prepare
the second vector freshly.  For every first label the old Gram and
the fresh Gram are both $JJ^*/s$.  The two final spectral readouts
used in Theorem~\ref{thm:rankspectrum92} attain the bounds.  All
three counted registers have dimension $s$.  This proves the assertion
also at $s=1$, where the antisymmetric receiver projection is zero.
\end{proof}
The three bounds apply at specified cuts.  They are not a cap on the
total workspace of an arbitrary implementation or a removal of the
reset architecture.  In particular, saturation remains the reset
value, not the unrestricted coherent-acquisition value.

\subsection{A stable common-projection theorem}
For density matrices $a,b$ with smaller rank at most $r$, set
\[
 \Delta_{d,r}(a,b)=d-r^{-1}
  -\|(\sqrt a\otimes\sqrt b)(I-dS)(\sqrt a\otimes\sqrt b)\|_1.
\]
When $(d,r)\ne(2,1)$ define
\begin{equation}\label{eq:rankrigidityconstant93}
 c_{d,r}=d-1-r^{-1}>0,\qquad
 K_{d,r}=r+\frac{(3+\sqrt2)^2}{c_{d,r}}.
\end{equation}
\begin{lemma}[Quantitative common support]\label{lem:projectionstability93}
For $d\ge2$, $1\le r\le d$ and $(d,r)\ne(2,1)$, there exists
a rank-$r$ projection $P$ such that
\begin{equation}\label{eq:projectionstability93}
 \max\{\|a-P/r\|_1^2,\|b-P/r\|_1^2\}
       \le K_{d,r}\Delta_{d,r}(a,b).
\end{equation}
Only the smaller Gram rank need be bounded by $r$; neither
invertibility nor equality of the two supports is assumed.
\end{lemma}
\begin{proof}
Put $f=\|\sqrt a\sqrt b\|_1$, $\eta=1-f^2$ and
$e=\tr(ab)-f^2/r$.  The overlap rank bound gives $e\ge0$, and
$f\le1$ gives $\eta\ge0$.  The proof of
Lemma~\ref{lem:rankswap92} gives
\begin{equation}\label{eq:rankdefects93}
              \Delta_{d,r}(a,b)\ge c_{d,r}\eta+e.
\end{equation}
Extend a right polar factor to a unitary $U$ so that
$Z=\sqrt a\sqrt bU=\sqrt{\sqrt a b\sqrt a}\succeq0$.
The Hilbert--Schmidt identity and the Schatten product inequality give
\[
 \|\sqrt bU-\sqrt a\|_2^2=2(1-f),\qquad
 \|Z-a\|_1\le\sqrt{2(1-f)}\le\sqrt{2\eta}.
\]
Choose any rank-$r$ projection $P$ containing the range of $Z$.
This is possible since $\operatorname{rank}Z\le r$.  Since
$\tr Z=f$, one has
\[
 \|Z-fP/r\|_2^2=e,\qquad
 \|Z-fP/r\|_1\le\sqrt{re}.
\]
These formulas remain valid at $f=0$, without normalizing $Z$.
As $1-f\le\sqrt\eta$, the triangle inequality yields
\[
       \|a-P/r\|_1\le(1+\sqrt2)\sqrt\eta+\sqrt{re}.
\]
The normalized coefficient vectors for $\sqrt a$ and $\sqrt bU$
have inner product $f$ and receiver marginals $a,b$.
Their pure-state trace distance is $2\sqrt{1-f^2}$, by its
rank-two spectral calculation.  Contractivity under partial trace
therefore gives $\|a-b\|_1\le2\sqrt\eta$, the trace-distance/root-fidelity estimate of
Fuchs--van de Graaf~\cite[Theorem~1]{FvdG93}.  Consequently
\[
       \|b-P/r\|_1\le(3+\sqrt2)\sqrt\eta+\sqrt{re}.
\]
Weighted Cauchy--Schwarz bounds the square of either right side by
$[r+(3+\sqrt2)^2/c_{d,r}](c_{d,r}\eta+e)$.
Now use \eqref{eq:rankdefects93}.  Every argument used matrices on
finite supports and a unitary extension of a partial isometry; no
inverse of a singular matrix or limit of an equality statement occurs.
\end{proof}
The constants in this lemma are explicit but not asserted to be optimal.
The projection may vary between branches; the theorem does not force
all receivers to occupy one common subspace throughout the protocol.

\begin{corollary}[Near-optimal projection decompositions]
\label{cor:nearoptimalframes93}
Assume $(d,r)\ne(2,1)$ and a strategy in
$\mathfrak R_{k,\ell}$ has reward $v\ge p_r-\varepsilon$.
Choose a dimension-preserving refined normal form; on its nonzero
branches put
\[
 w_{yh}=\tr A_{yh},\qquad a_{yh}=A_{yh}/w_{yh},\qquad
 \mu_{yh}=w_{yh}/d,\qquad
 B_\varepsilon=\frac{2K_{d,r}d(d+1)}{t^2}\varepsilon.
\]
There are rank-$r$ projections $P_{yh}$ such that
\begin{align}
 \sum_{y,h}\mu_{yh}
  \bigl(\|a_{yh}-P_{yh}/r\|_1^2+
              \|b_{yh}-P_{yh}/r\|_1^2\bigr)&\le2B_\varepsilon,
 \label{eq:nearoptimalframes93}\\
 \frac1d\sum_y\left\|\rho-\sum_hw_{yh}P_{yh}/r\right\|_1^2
          &\le B_\varepsilon.
 \label{eq:framebarycenter93}
\end{align}
Finite and countable recorded instruments are allowed.  The weights
$\mu_{yh}$ sum to one and are Gram weights, not probabilities under a
chosen hypothesis.  For public mixtures the same statements hold with
the public seed adjoined to the refined record, using the common
purified initial Gram matrix.
\end{corollary}
\begin{proof}
Subtract \eqref{eq:unsaturatedrankupper93} from the attained class
value.  Since $\sum_{y,h}w_{yh}=d$, this gives
\[
 \sum_{y,h}\mu_{yh}\Delta_{d,r}(a_{yh},b_{yh})
       \le\frac{\varepsilon}{\alpha d}
       =\frac{2d(d+1)}{t^2}\varepsilon.
\]
Apply Lemma~\ref{lem:projectionstability93} to each branch and sum.
For the second estimate, each fixed $y$ has $\sum_hw_{yh}=1$ and
$\sum_hw_{yh}a_{yh}=\rho$.  The triangle inequality followed by
Jensen's inequality bounds its squared residual by
$\sum_hw_{yh}\|a_{yh}-P_{yh}/r\|_1^2$.
Average over $y$.  The proof uses nonnegative sums with finite total
weight, so countable traces pass to the limit.  The seed construction
in Lemma~\ref{lem:resetnormal91}, followed by the dimension-preserving
refinement, preserves the common Gram sum and the achieved reward.
\end{proof}
The total $\mu$-weight on branches for which either deviation is at
least $u>0$ is at most $2B_\varepsilon/u^2$.  This is a quantitative
statement at the relative loss scale $\varepsilon/t^2$, not a claim
of vanishing error at fixed nonzero calibration defect.

\subsection{A spectral loss certificate}
Let
\[
 \mathcal C_{d,r}=\{\sigma\in\mathcal D_d:\sigma\preceq I/r\},
 \qquad e_r(\rho)=\tr(\rho-I/r)_+.
\]
\begin{proposition}[Distance to the optimal initial spectra]
\label{prop:spectralloss93}
For every density matrix $\rho$,
\begin{equation}\label{eq:capdistance93}
 \inf_{\sigma\in\mathcal C_{d,r}}\|\rho-\sigma\|_1=2e_r(\rho).
\end{equation}
For $(d,r)\ne(2,1)$, any strategy in $\mathfrak R_{k,\ell}$
with initial Gram $\rho$ and achieved equal-prior reward $v$ satisfies
\begin{equation}\label{eq:spectralloss93}
 p_r-v\ge \frac{2t^2}{K_{d,r}d(d+1)}\,e_r(\rho)^2.
\end{equation}
\end{proposition}
\begin{proof}
Let $Q$ be the spectral projection of $\rho$ above $1/r$.
For $\sigma\preceq I/r$, one has
$\tr Q(\rho-\sigma)\ge e_r(\rho)$.
As $\tr(\rho-\sigma)=0$, positive-part optimization gives
$\|\rho-\sigma\|_1\ge2e_r(\rho)$.
Conversely, in an eigenbasis of $\rho$, cap every eigenvalue exceeding
$1/r$ and redistribute the removed mass among eigenvalues below $1/r$.
There is enough capacity since $d/r\ge1$.  This produces a density
matrix in $\mathcal C_{d,r}$ at trace distance exactly $2e_r(\rho)$.

Every matrix $\sum_hw_{yh}P_{yh}/r$ in
\eqref{eq:framebarycenter93} belongs to $\mathcal C_{d,r}$.
Set $\varepsilon=p_r-v\ge0$ there and use
\eqref{eq:capdistance93}; rearranging $4e_r(\rho)^2\le
2K_{d,r}d(d+1)(p_r-v)/t^2$ proves the bound.
\end{proof}
The estimate concerns the actual first-input marginal, up to its fixed
transpose, and a specified ideal reset architecture.  Under the
calibrated perturbations of Corollary~\ref{cor:gamestability91}, the
same strategy's nominal reward differs from its observed reward by at
most the stated channel, prior and latent-law error; a preparation
defect adds its stated half-trace contribution.  That total error
must be added to $p_r-v$ before applying the certificate.  There is no
uncalibrated inference of the reset cut from a score.

\subsection{The exceptional qubit face}
\begin{proposition}[Pinching at the exceptional endpoint]
\label{prop:qubitpinching93}
Let $d=2$, $r=1$, and suppose $a$ is pure.  Put
$\mathcal P_a(b)=aba+(I-a)b(I-a)$ and
$\Delta=\Delta_{2,1}(a,b)$.  Then
\begin{equation}\label{eq:qubitpinching93}
 \|b-\mathcal P_a(b)\|_1^2=2\Delta-\Delta^2\le2\Delta.
\end{equation}
In particular $\Delta=0$ if and only if $[a,b]=0$.
If instead $b$ is the pure factor, interchange $a,b$.
\end{proposition}
\begin{proof}
By simultaneous unitary conjugation take $a=|0\rangle\langle0|$ and
$b=\left(\begin{smallmatrix}x&z\\\bar z&1-x\end{smallmatrix}\right)$.
The nonzero block of the centered filtered swap is
$\sqrt b(I-2a)\sqrt b$.  Its trace is $1-2x$ and its determinant
is $-\det b$.  Its trace norm, also when $b$ is singular, is therefore
\[
 \sqrt{(1-2x)^2+4\det b}=\sqrt{1-4|z|^2}=1-\Delta.
\]
The pinching difference has eigenvalues $\pm|z|$.
Squaring gives \eqref{eq:qubitpinching93}.  Vanishing is exactly
$z=0$, or $[a,b]=0$.
\end{proof}
Thus the exceptional equality set includes a pure factor with any
commuting mixed qubit state, not just identical or orthogonal pure
factors.  A common rank-one projection cannot replace pinching here.
The initial-spectrum theorem remains valid because
$\mathcal C_{2,1}=\mathcal D_2$.  This exception is separate from the
zero-signal case $t=0$, in which every strategy has reward $1/2$.

\section{Exact values at a prescribed initial spectrum}
\label{sec:spectralvalue94}
The equality set of a resource bound does not determine the value away
from that set.  We now solve this remaining optimization for the biased
basis experiment of Theorem~\ref{thm:operationalhierarchy90}.  Besides
giving the exact spectral loss, the solution distinguishes protocols
that merely retain a receiver from protocols that communicate a
receiver outcome to the next preparation.

Fix $d\ge2$, a weighted exact second-moment family satisfying
\eqref{eq:basismoments90}, and $0\le t\le1$.  The alternative basis is
selected once and used at both calls.  In this section only, the priors
are those of \eqref{eq:gameprior90}; write
\[
 L_0=2(d+1)-t^2,\qquad p_E=(d+1)/L_0.
\]
Fix the \emph{pure} first acquisition through its Gram matrix
$\rho=C_0^*C_0\in\mathcal D_d$.  Its input marginal is
$\rho^{\mathsf T}$.  The initial acquisition is not replaced by a
classically resolved mixture.  Subsequent instruments, mixtures of
fresh states, public seeds and tester limits retain the conventions of
Definition~\ref{def:rankclass92}.

For a nonnegative vector $u$, put $v_i=\sqrt{u_i}$ and define
\begin{equation}\label{eq:secularmatrix94}
 \chi(u)=\lambda_{\max}\bigl(vv^*-\diag u\bigr).
\end{equation}
Appending zeros does not change this nonnegative number.  In the secular
equation omit zero coordinates; the matrix definition remains valid on
every support face.  If at least
two entries of $u$ are positive, it is the unique positive solution of
\begin{equation}\label{eq:secularroot94}
                 \sum_i\frac{u_i}{u_i+\chi}=1.
\end{equation}
If at most one entry is positive, set $\chi(u)=0$.

Let $\lambda_1\ge\cdots\ge\lambda_d\ge0$ be the eigenvalues of
$\rho$.  For $1\le r\le d$, choose the smallest
$m\in\{0,\ldots,r-1\}$ for which
\[
 \lambda_{m+1}\le\tau,\qquad
 \tau=\frac{\sum_{i>m}\lambda_i}{r-m},
\]
and define the rank-$r$ majorant
\begin{equation}\label{eq:rankmajorant94}
 q^{(r)}(\lambda)
   =(\lambda_1,\ldots,\lambda_m,
               \underbrace{\tau,\ldots,\tau}_{r-m},0,\ldots,0).
\end{equation}
The choice always exists at $m=r-1$.  The preceding failed test, if
$m>0$, implies $\lambda_m>\tau$, so the displayed vector is decreasing.
Zero tails are allowed; ``rank-$r$'' here means rank at most $r$.
Write $\mathcal H_r(\rho)=\chi(q^{(r)}(\lambda(\rho)))$.
The ``smallest $m$'' convention fixes the active set at a water-level
tie.  Entries equal to $\tau$ belong to the flattened tail.  Including
such a tied entry in the initial segment would give the same vector,
so $q^{(r)}$ is unique even there.  If $r$ reaches the support size,
$q^{(r)}=\lambda$, padded by zeros.  These are the conventions for
$\chi$ and $q^{(r)}$ throughout the article; eigenbasis choices are not
part of their definitions.  Lemma~\ref{lem:spectralroof94} proves the
least-majorant assertion and hence uniqueness intrinsically.

\begin{lemma}[Continuity at water-level ties]\label{lem:watermap96}
On decreasing probability vectors the map $q^{(r)}$ is continuous,
piecewise affine, and nonexpansive in $\ell^1$:
\[
 \|q^{(r)}(\lambda)-q^{(r)}(\mu)\|_1
                                     \le\|\lambda-\mu\|_1.
\]
\end{lemma}
\begin{proof}
For each fixed active index $m$, the displayed rule is an affine map:
its first $m$ coordinates are unchanged, and each remaining input
coordinate contributes $1/(r-m)$ to each of the next $r-m$ outputs.
Its matrix has nonnegative entries and every column sums to one, so
its induced $\ell^1$ norm is one.  The finitely many closed active
regions are polyhedral.  On their overlaps the water-level equalities
make the formulas identical, including zero tails.  Hence the map is
continuous.  A line segment in the decreasing simplex has a finite
partition into these regions; summing the norm bounds along its
collinear subsegments proves the assertion.
\end{proof}

\begin{theorem}[Prescribed-spectrum value]\label{thm:spectralvalue94}
For the once-selected finite family and the fixed pure first acquisition
above, the optimum over the complete class $\mathfrak R_{k,\ell}$,
$1\le k,\ell\le d$, is
\begin{equation}\label{eq:spectralvalue94}
 P_{k,\ell}^{\rm b}(\rho)
       =p_E+\frac{t^2}{2L_0}\mathcal H_r(\rho),
                 \qquad r=\min\{k,\ell\}.
\end{equation}
At $r=1$ the spectral gain is zero.  At $t=0$ the value is $1/2$.
If $r\ge\operatorname{rank}\rho$, the majorant is the initial spectrum
itself.  Trace norms used in the derivation are unhalved; the factor
$1/2$ in \eqref{eq:spectralvalue94} converts the filtered-swap leaf
normalization into this Bayes gain.
The optimum is attained by an instrument with at most $d$ recorded
outcomes, the same instrument after every first device label, and fresh
references of dimension at most $r$.  Both retained quantum registers
may have dimension $r$.  The second preparation may depend on the
recorded receiver outcome, but need not depend on the first device
label.  This construction acts on the prescribed initial acquisition;
it does not optimize over its Gram matrix.
\end{theorem}
For $r=d$ no instrument is needed and the answer is given directly by
the secular equation for $\lambda(\rho)$.  The theorem gives an exact
value for the specified priors.  The equal-prior spectral rigidity of
Section~\ref{sec:spectralrigidity93} remains a separate result.

\subsection{The scalar function and a rank-constrained leaf}
\begin{lemma}[Secular concavity]\label{lem:secular94}
The function $\chi$ is continuous, symmetric, homogeneous of degree one,
coordinatewise nondecreasing and concave on each nonnegative orthant.
It is strictly Schur-concave on each fixed positive trace simplex:
if $u\prec w$ and $u,w$ are not permutations, then
$\chi(u)>\chi(w)$.  Here vectors are padded to a common length.
\end{lemma}
\begin{proof}
On a support containing at least two coordinates, the matrix in
\eqref{eq:secularmatrix94} has strictly positive off-diagonal entries.
A maximizing Rayleigh vector can be chosen nonnegative, and its
coordinates on this support must be positive.  Its eigenvalue equation
is $(u_i+\chi)x_i=\sqrt{u_i}\langle v,x\rangle$.
The Rayleigh value is positive, and summing gives
\eqref{eq:secularroot94}.  Its left side strictly decreases from the
support size to zero, proving uniqueness.  The support-one case follows
directly from the matrix formula.  That formula also gives continuity,
symmetry and homogeneity, including at the boundary.

For $a>0$, the root characterization says
\[
 \chi(u)\ge a\quad\Longleftrightarrow\quad
                     \sum_i\frac{u_i}{u_i+a}\ge1.
\]
Each summand is increasing and concave in $u_i$, so every such
superlevel set is convex and coordinatewise increasing.  If
$c=\chi(u)>0$ and $e=\chi(w)>0$, convexity of the unit superlevel set
applied with weights $c/(c+e),e/(c+e)$ gives
$\chi(u+w)\ge c+e$.  If one value vanishes the same inequality follows
from coordinate monotonicity.  Homogeneity now proves concavity.

Symmetry and concavity imply Schur concavity by convex combinations of
permutations.  For strictness, if $\chi(w)>0$, strict concavity of
$x/(x+\chi(w))$ shows that $u\prec w$, without permutation equality,
implies
$\sum_i u_i/(u_i+\chi(w))>1$.  The root for $u$ is therefore larger.
If $\chi(w)=0$, $w$ has only one positive entry and any nonpermutation
vector majorized by it has at least two, again giving strictness.
\end{proof}

For $A,B\succeq0$ let
\[
 \mathcal N(A,B)=\tr\bigl[-(\sqrt A\otimes\sqrt B)S
                                  (\sqrt A\otimes\sqrt B)\bigr]_+.
\]
\begin{lemma}[Optimal fresh Gram at a fixed old Gram]
\label{lem:spectralleaf94}
If $a_1\ge\cdots\ge a_d\ge0$ are the eigenvalues of $A$, then
\begin{equation}\label{eq:spectralleaf94}
 \max_{B\in\mathcal D_d,\ \operatorname{rank}B\le\ell}
                  \mathcal N(A,B)
           =\frac12\chi(a_1,\ldots,a_\ell).
\end{equation}
A maximizing $B$ commutes with $A$ and is supported on its leading
$\ell$ eigenspaces.  If $c=\chi(a_1,\ldots,a_\ell)>0$, an explicit
choice on that support is
\begin{equation}\label{eq:optimalfresh94}
 b_i=\frac{a_i/(a_i+c)^2}{\sum_{j\le\ell}a_j/(a_j+c)^2}.
\end{equation}
No invertibility of $A$ is required.  A commuting optimizer exists;
if a leading eigenspace is degenerate the chosen eigenbasis, and hence
the displayed optimizer, need not be unique.  The factor $1/2$ in
\eqref{eq:spectralleaf94} concerns negative mass, with no halving of
the trace norm in Lemma~\ref{lem:swapfilter90}.
\end{lemma}
\begin{proof}
Lemma~\ref{lem:swapfilter90} expresses $2\mathcal N(A,B)$ as
$(\tr\sqrt Z)^2-\tr Z$, where $Z=\sqrt A B\sqrt A$.
First work on the support of $A$.  Compressing $B$ to that support
leaves $Z$ unchanged and does not increase its rank.  If the compressed
trace is positive, normalizing it can only increase the nonnegative,
homogeneous objective.  If the trace vanishes the objective is zero.
We may therefore take $A$ invertible on its support and
$\tr(A^{-1}Z)=1$.

For fixed eigenvalues $z_1\ge\cdots\ge z_s>0$ of $Z$, $s\le\ell$,
rearrangement gives
\[
                  \sum_{i=1}^s z_i/a_i\le\tr(A^{-1}Z)=1.
\]
For completeness, in eigenbases the trace is a sum of products weighted
by squared entries of a unitary, a doubly stochastic matrix.  Its
minimum over that larger convex set is achieved at a permutation;
exchanging two inverted pairs puts the largest $z_i$ on the largest
$a_i$.  This proves the inequality.  Put $Z'$ diagonal with these
$z_i$ on the leading eigenspaces of $A$ and
$B'=A^{-1/2}Z'A^{-1/2}$.  Then $0<\tr B'\le1$;
normalizing $B'$ proves that a commuting maximizer with the asserted
support suffices.

For such a $B$, write $x_i=\sqrt{b_i}$, so $\|x\|_2=1$.  The objective
is
\[
 2\mathcal N(A,B)
       =\left(\sum_{i\le\ell}\sqrt{a_i}x_i\right)^2
                          -\sum_{i\le\ell}a_ix_i^2.
\]
The Rayleigh maximum is \eqref{eq:secularmatrix94}, with a nonnegative
maximizer.  Its eigenvector equation yields
\eqref{eq:optimalfresh94}.  With at most one positive $a_i$ the maximum
is zero, achieved by any rank-one Gram on the support, or any state if
$A=0$.  This handles every singular case explicitly.
\end{proof}

\subsection{A spectral evaluation of the concave envelope}
The next lemma records the majorization step separately.  It applies
to every continuous symmetric concave function, not only to the leaf
objective just derived.  The underlying majorization and ensemble
methods are classical; their operational use in entanglement
transformation and assistance is discussed in
\cite{Nielsen94,JonathanPlenio94,Assistance94}.

\begin{lemma}[Rank-capped spectral envelopes]\label{lem:spectralroof94}
Let $f$ be a continuous symmetric concave function on the probability
simplex of length $r$.  For a density matrix of rank at most $r$, pad
its eigenvalue vector to length $r$ before applying $f$.  Then
\begin{equation}\label{eq:spectralroof94}
 \max_{\substack{\rho=\sum_hw_ha_h\,,\ w_h\ge0\\
                         a_h\in\mathcal D_d^{\le r}}}
        \sum_hw_h f(\lambda(a_h))
        =f(q^{(r)}_1(\lambda(\rho)),\ldots,
                                   q^{(r)}_r(\lambda(\rho))).
\end{equation}
The maximum admits at most $d$ atoms commuting with $\rho$, all with
the same spectrum $q^{(r)}$.  In particular $\mathcal H_r$ is concave
on $\mathcal D_d$; its homogeneous extension to positive matrices is
concave and superadditive.
\end{lemma}
\begin{proof}
Write $q=q^{(r)}(\lambda)$.  Its construction gives $\lambda\prec q$.
Indeed its initial $m$ partial sums agree with those of $\lambda$;
each subsequent entry of $\lambda$ is at most $\tau$, and its total
sum is one.  Moreover $q$ is majorized by every decreasing probability
vector $z$ of rank at most $r$ which majorizes $\lambda$.  Up to $m$
this follows from the majorization inequalities.  Between $m$ and $r$
the decreasing entries of $z$ give
\[
 \sum_{i\le j}z_i\ge
 \frac{r-j}{r-m}\sum_{i\le m}z_i+\frac{j-m}{r-m}
 \ge\sum_{i\le m}\lambda_i+(j-m)\tau
                   =\sum_{i\le j}q_i.
\]
After $r$ both totals are one.  Thus $q$ is the least such majorant.

For any finite decomposition, the variational formula for the sum of
the $j$ largest eigenvalues gives
\[
 \lambda(\rho)\prec z:=\sum_hw_h\lambda(a_h),
\]
where the eigenvalue vectors are decreasing and padded to length $d$.
The vector $z$ has rank at most $r$.  Concavity and Schur concavity give
$\sum_hw_hf(\lambda(a_h))\le f(z)\le f(q)$.

For the converse, $\lambda\prec q$ implies that $\lambda$ lies in the
convex hull of the permutations of $q$.  One direct verification is
as follows.  For every real vector $b$, rearrangement followed by
summation by parts in the partial-sum inequalities gives
$b\cdot\lambda\le b^\downarrow\cdot\lambda
 \le b^\downarrow\cdot q=\max_\pi b\cdot\pi q$.
A separating hyperplane is therefore impossible.  This convex hull
lies in a hyperplane of dimension at most $d-1$, so Carath\'eodory
reduces the combination to at most $d$ permutations.  In an eigenbasis
of $\rho$ they give commuting density matrices $a_h$ of spectrum $q$.
Their $f$ values coincide, proving attainment.

Concavity follows also by combining any two attaining decompositions.
For a countable decomposition, the same upper follows from the Ky Fan
inequalities and Jensen's inequality by convergence of its bounded
trace weights.  At zero trace define the homogeneous extension to be
zero.  Trace-weighted convex combinations prove its concavity and
superadditivity.  Apply this with $f=\chi$ for the final assertion.
\end{proof}

For clarity, the homogeneous extension just proved is
\[
 \widehat{\mathcal H}_r(A)=
 \begin{cases}
  (\tr A)\mathcal H_r(A/\tr A),&\tr A>0,\\
  0,&A=0.
 \end{cases}
\]
For positive $A,B$ of traces $a,b>0$, concavity on densities gives
$\widehat{\mathcal H}_r(A+B)\ge
\widehat{\mathcal H}_r(A)+\widehat{\mathcal H}_r(B)$ by weights
$a/(a+b)$ and $b/(a+b)$; zero traces are immediate.  Subsequent formulas
write $\mathcal H_r(A)$ for this extension on positive matrices.

\subsection{Proof of the prescribed-spectrum theorem}
\begin{proof}[Proof of Theorem~\ref{thm:spectralvalue94}]
Use the dimension-preserving normal form of
Lemma~\ref{lem:ranknormal92}.  The fixed initial Gram gives
$\sum_hA_{yh}=\rho$ for every $y$.  Put $w_{yh}=\tr A_{yh}$ and
$a_{yh}=A_{yh}/w_{yh}$ on nonzero branches.  These are Gram weights;
they are not substituted for the hypothesis-dependent history
probabilities.  Mixed fresh states and receiver channels are refined
classically, without increasing either retained quantum dimension.

The biased payoff blocks are $-\kappa S$ for equal labels and
$-2\kappa\Pi_-/(d-1)\preceq0$ for unequal labels, where
$\kappa=t^2/(dL_0)$.  Thus only equal labels can improve on the
constant alternative decision.  Lemma~\ref{lem:spectralleaf94} gives
\[
 P-p_E\le\frac{\kappa}{2}\sum_{y,h}
                  \chi(\lambda_1(A_{yh}),\ldots,\lambda_\ell(A_{yh})).
\]
Each summand is at most $\mathcal H_r(A_{yh})$.  If $k\le\ell$, the
rank of $A_{yh}$ is at most $r=k$ and the two expressions agree.  If
$\ell<k$, the first $r=\ell$ coordinates of its rank-$r$ majorant
are coordinatewise at least its leading $r$ eigenvalues; use the
monotonicity of $\chi$.  Superadditivity then gives
\[
 P-p_E\le\frac{\kappa}{2}\sum_y\mathcal H_r(\rho)
                    =\frac{t^2}{2L_0}\mathcal H_r(\rho).
\]
The same bound survives public mixing, countable normal instruments
and tester closure, by the fixed-Gram normal form and trace tails.
At $t=0$ all devices coincide and both sides reduce to $1/2$.

For attainment take the commuting decomposition
$\rho=\sum_hw_ha_h$ supplied by Lemma~\ref{lem:spectralroof94}, all
$a_h$ with spectrum $q^{(r)}$.  On each support choose rectangular
$C_h,D_h:\C^d\to\C^r$ with
\[
 C_h^*C_h=w_ha_h,\qquad D_h^*D_h=b_h,
\]
where $b_h$ is the permuted optimizing fresh Gram in
\eqref{eq:optimalfresh94}.  If $\chi(q^{(r)})=0$, use any pure
fresh Gram and the constant decision.  Otherwise $b_h$ is a density
matrix of rank at most $r$, supported with $a_h$.  The maps
$V_h=C_hC_0^+$ satisfy $V_hC_0=C_h$ and
$\sum_hV_h^*V_h=I$ on the initial Schmidt support: positivity of the
Gram sum implies that every $C_h$ vanishes on $\ker\rho$.
They therefore form an instrument on the given initial receiver.
Proposition~\ref{prop:finiteinstrument96} gives its domains, codomains,
completion on an unused orthogonal complement and the at-most-$d$
message bound explicitly, without retaining an extra environment.

Use this same instrument after every $y$ and prepare $|D_h\rangle$
independently conditional on its recorded outcome.  For equal labels,
the positive spectral projection of the filtered negative swap attains
$\kappa w_h\chi(q^{(r)})/2$.  For unequal labels choose the zero
scalar-hypothesis effect.  Complements complete each binary readout.
Summing the $d$ first labels gives equality.  These unnormalized
formulas also define the endpoint $t=1$ and null histories.  The
construction proves attainment without an extra retained environment.
\end{proof}

\subsection{When a receiver message is strictly useful}
The following comparison fixes the initial acquisition and changes
only the permission to send a receiver outcome to the fresh source.
Let $P_{d,r}^{\rm nm}(\rho)$ be the optimum with old receiver dimension
$d$ and fresh reference dimension $r$, when the fresh state may depend
on the first device label and an independent public seed, but on no
outcome obtained by measuring the old receiver.  Arbitrary receiver
processing and the final joint readout remain allowed.  Any old
processing without such communication can be postponed to the final
readout by Lemma~\ref{lem:postponement96}; the old bound $d$ contains
the entire initial Schmidt support.  This is a subclass of $\mathfrak R_{d,r}$, not the
complete-measure-and-reprepare class.

\begin{corollary}[Exact value of a receiver message]\label{cor:messagevalue94}
For the fixed experiment and initial Gram above,
\begin{equation}\label{eq:nomessage94}
 P_{d,r}^{\rm nm}(\rho)
          =p_E+\frac{t^2}{2L_0}\chi(\lambda_1,\ldots,\lambda_r).
\end{equation}
If $t>0$, then
\begin{equation}\label{eq:messagestrict94}
 P_{d,r}^{\rm b}(\rho)>P_{d,r}^{\rm nm}(\rho)
                 \quad\Longleftrightarrow\quad
                      2\le r<\operatorname{rank}\rho.
\end{equation}
The strict gain is realized with at most $d$ receiver outcomes and
without first-device-label dependence of the fresh state.
\end{corollary}
\begin{proof}
Without a receiver message each first-label branch retains Gram
$\rho$.  Lemma~\ref{lem:spectralleaf94} gives the displayed upper,
and the commuting optimizer with the positive-part final decisions
attains it.  Public mixtures cannot improve an already maximized
linear score.  If $r=1$, both secular values vanish.  If
$r\ge\operatorname{rank}\rho$, no tail is moved and $q^{(r)}=\lambda$.
In the remaining case the leading $r$ eigenvalues are all positive,
and $q^{(r)}$ increases at least one of them without decreasing any.
For a vector with at least two positive entries the root strictly
increases when any positive coordinate is increased, as follows
immediately from \eqref{eq:secularroot94}.  This proves both directions.
\end{proof}

For an explicit example take $d=3$, $r=2$ and
$\rho=\diag(3/4,1/8,1/8)$.  Then
\[
 q^{(2)}=(3/4,1/4,0),\quad
 \chi(q^{(2)})=\sqrt3/4,\quad
 \chi(\lambda_1,\lambda_2)=\sqrt6/8.
\]
With $L_0=8-t^2$, the exact advantage of the message is
\begin{equation}\label{eq:messageexample94}
 \frac{t^2(2\sqrt3-\sqrt6)}{16L_0}>0\qquad(t>0).
\end{equation}
Use two Gram atoms
$a_1=\diag(3/4,1/4,0)$ and $a_2=\diag(3/4,0,1/4)$ with
weights $1/2$.  With $C_0=\sqrt\rho$, explicit old-receiver maps are
\[
 V_1=\begin{pmatrix}1/\sqrt2&0&0\\0&1&0\end{pmatrix},\qquad
 V_2=\begin{pmatrix}1/\sqrt2&0&0\\0&0&1\end{pmatrix}.
\]
They satisfy $\sum_hV_h^*V_h=I_3$ and
$C_h=V_hC_0$, $C_h^*C_h=a_h/2$.
The fresh Grams are respectively
$\diag(1/2,1/2,0)$ and $\diag(1/2,0,1/2)$; normalized fresh
vectors have the corresponding two equal Schmidt weights.  The final
positive-part readouts used in the proof attain the claimed scores.
Thus the improvement is a realizable recorded instrument, not a
change of initial state or postselection on a favorable history.

\begin{corollary}[Spectral loss and exact saturation]\label{cor:spectralsaturation94}
For $r>1$, the difference from the optimized rank-$r$ value of
Theorem~\ref{thm:rankspectrum92} is exactly
\begin{equation}\label{eq:exactspectralloss94}
 P_{k,\ell}^{\rm b}-P_{k,\ell}^{\rm b}(\rho)
       =\frac{t^2}{2L_0}\left(\frac{r-1}{r}-\mathcal H_r(\rho)\right).
\end{equation}
At $t>0$ it vanishes exactly when $\rho\preceq I/r$.
For $t>0$ and a fixed $\rho$ of rank $s$, its constrained optimum strictly
increases as the smaller receiver dimension increases from $1$ to
$s$, and is constant thereafter.  The smallest dimension attaining
its full-reset value is therefore $s$, including the trivial case
$s=1$.  At $t=0$ every dimension gives the same value $1/2$.
\end{corollary}
\begin{proof}
The uniform vector of length $r$ is majorized by every probability
vector of that length.  Lemma~\ref{lem:secular94} gives
$\chi(q)\le(r-1)/r$, with equality only at that uniform vector.
The construction of $q^{(r)}$ makes this equivalent to
$\lambda_1\le1/r$.  Subtracting \eqref{eq:spectralvalue94} from
\eqref{eq:rankbiased92} proves the exact loss identity.

The least-majorant property gives
$q^{(r+1)}\prec q^{(r)}$.  If $r<s$, these have respectively $r+1$
and $r$ positive entries and are not permutations, so strict Schur
concavity gives a strict increase.  At $r\ge s$ both vectors equal
$\lambda(\rho)$.  This proves the saturation statement.
\end{proof}

The secular root and the majorant are functions of a supplied spectrum.
Bisection of \eqref{eq:secularroot94} evaluates the value without an
optimization over physical protocols; a finite convex decomposition
then realizes it.  Real input spectra and real weights are permitted.
These statements do not assume a rational eigenbasis or give an
efficient construction of the unknown-device design.  In particular,
the general gain from a receiver message proved here is compatible
with the no-feedback globally attaining preparations in
Corollary~\ref{cor:threecuts93}: those preparations optimize the initial
Gram, whereas \eqref{eq:messagestrict94} holds it fixed.

\section{Finite attainment on a prescribed receiver}
\label{sec:attainment96}
We make the spectral converse operational on the actual initial
Schmidt support.  This also records all normalization factors for a
three-dimensional example.  Neither construction replaces a fixed
pure initial acquisition by an announced mixture.

\begin{proposition}[Support-exact finite instrument]
\label{prop:finiteinstrument96}
Let $C_0:\C^d\to R$ satisfy $C_0^*C_0=\rho$, $\tr\rho=1$, and
let $S_0=\operatorname{ran}C_0\subseteq R$.  Suppose
$\rho=\sum_{h=1}^n w_h a_h$, where $w_h>0$, $\sum_hw_h=1$ and
$\operatorname{rank}a_h\le r$.  For each $h$ let $b_h$ be a density
matrix of rank at most $\ell$.  Choose coefficient maps
\[
 C_h:\C^d\to\C^r,\quad C_h^*C_h=w_ha_h,
 \qquad D_h:\C^d\to\C^\ell,\quad D_h^*D_h=b_h.
\]
The Moore--Penrose inverse $C_0^+:R\to\C^d$, zero on $S_0^\perp$,
gives $V_h=C_hC_0^+:R\to\C^r$ with
\begin{equation}\label{eq:finiteinstrument96}
 V_hC_0=C_h,\qquad \sum_h V_h^*V_h=P_{S_0}.
\end{equation}
These maps extend to an instrument on $R$ with the same $n$ classical
outcomes and quantum outputs of dimension at most $r$.  With the fresh
vector $|D_h\rangle$ prepared independently after outcome $h$, the
receiver block is the product in \eqref{eq:resetblock91}.
For the spectral decomposition used in
Theorem~\ref{thm:spectralvalue94}, $n\le d$ suffices.
\end{proposition}
\begin{proof}
If $x\in\ker\rho$, positivity and the barycenter equation imply
$\langle x,a_hx\rangle=0$ for every $h$.  Hence $C_hx=0$ and
$C_hP_{\operatorname{supp}\rho}=C_h$.  The Moore--Penrose identities
then give $V_hC_0=C_h$ and
\[
 \sum_hV_h^*V_h=(C_0^+)^*\rho C_0^+=P_{S_0}.
\]
This proves \eqref{eq:finiteinstrument96} even when $\rho$ is singular.
To complete the instrument without adding a recorded outcome, choose
a unit vector $e\in\C^r$ and add to its first branch the completely
positive map
\[
             X\longmapsto\tr(P_{S_0^\perp}X)|e\rangle\langle e|.
\]
Together with $X\mapsto V_1XV_1^*$ in that branch and
$X\mapsto V_hXV_h^*$ in the others, the total map is trace
preserving.  This completion vanishes on all physically occupied
initial receiver states.  It is normal also on a larger separable
receiver space, and retains no dilation environment.  The fresh vector
has squared norm $\tr b_h=1$.  Measurement of each input gives the
product block by expansion in the fixed input bases.  The final POVMs
therefore implement the chosen leaf optima.

For the spectral theorem, put $q=q^{(r)}(\lambda(\rho))$.
The least-majorant proof in Lemma~\ref{lem:spectralroof94} gives
$\lambda(\rho)\prec q$.  A decreasing vector majorized by $q$ lies
in the convex hull of its permutations; this follows by successive
two-coordinate averaging, or equivalently by the doubly stochastic
characterization of majorization.  Every permutation belongs to the
affine hyperplane $\sum_i x_i=1$, of dimension $d-1$.  Carath\'eodory
therefore reduces its convex combination to at most $d$ terms.
Use those permutations as diagonal atoms in a fixed eigenbasis of
$\rho$ and apply the construction above.  Eigenvalue ties can change
the chosen eigenbasis or decomposition, not the value or the bound.
The weights may be real; this argument asserts no bit or gate
complexity bound.
\end{proof}
The probability of $(y,h)$ under a first measurement $E^\theta$ is
$w_h\tr[a_h(E_y^\theta)^{\mathsf T}]$.  Thus $w_h$ is generally not
the conditional receiver-outcome probability after observing $y$.
No physical branch is postselected in the attaining construction.

\subsection{One complete three-dimensional calculation}
Let $d=3$, $r=2$, $t=1$, and
\[
 \rho=\diag(3/4,1/8,1/8),\qquad
 |C_0\rangle=\frac{\sqrt3}{2}|1,1\rangle+
            \frac{|2,2\rangle+|3,3\rangle}{2\sqrt2}.
\]
Its squared norm is one.  For a fully specified finite device family,
put $\omega=e^{2\pi i/3}$.  Use the coordinate basis and the three
bases
\[
 |u_j^{(a)}\rangle=\frac1{\sqrt3}
       \sum_{x=0}^2\omega^{a x^2+jx}|x+1\rangle,
       \qquad a=0,1,2,\quad j=0,1,2,
\]
with all six orders of each basis.  These $24$ ordered bases have
weights $1/24$.  Their projections satisfy
\eqref{eq:basismoments90}.  To check the equal-label identity directly,
consider the $12$ vectors before ordering.  For their nine chirp
vectors the $(x,u;x',u')$ tensor entry in the sum of projectors
squared is
\[
 \frac19\sum_{a,j=0}^2
   \omega^{a(x^2+u^2-x'^2-u'^2)+j(x+u-x'-u')}.
\]
It equals one exactly when the two multisets $\{x,u\}$ and
$\{x',u'\}$ agree, and zero otherwise: equal sums and square sums
over $\mathbb F_3$ imply equal products.  The three coordinate
vectors add one to the fully diagonal entries.  The full sum is
therefore $I+S$.  Dividing by $12$ proves the equal-label moment;
summing the other two labels gives $(3I-S)/24$ for distinct labels.
This supplies the exact weighted family, rather than only its existence.

Under the scalar hypothesis $C_y=I/3$.  Under the alternative one of
the $24$ ordered projective measurements is drawn \emph{once} and
used twice.  The priors are $p_C=3/7$, $p_E=4/7$, so each individual
alternative has unconditional prior $1/42$.  All these weights are
rational, whereas the vectors and values below are algebraic.

\begin{proposition}[A prescribed-spectrum message witness]
\label{prop:qutritwitness96}
For this experiment and initial acquisition the exact communicated
and no-message values with old dimension three and fresh dimension
two are
\begin{equation}\label{eq:qutritvalues96}
 P_{3,2}^{\rm b}(\rho)=\frac47+\frac{\sqrt3}{56},\qquad
 P_{3,2}^{\rm nm}(\rho)=\frac47+\frac{\sqrt6}{112}.
\end{equation}
Two receiver outcomes attain the first value using an old output of
dimension two.  The second value retains the complete old receiver
and is an upper over all no-message strategies, not just the exhibited
one.  The strict gain is $(2\sqrt3-\sqrt6)/112$.
\end{proposition}
\begin{proof}
The water-filling rule has $m=1$, $\tau=1/4$, hence
$q^{(2)}=(3/4,1/4,0)$.  Its secular root is $\sqrt3/4$, whereas
the leading two original eigenvalues have root $\sqrt6/8$.
For two positive entries $(a,b)$ the secular equation reduces to
$x^2=ab$.  In particular the disjoint rational enclosures
\[
 \frac{433}{1000}<\frac{\sqrt3}{4}<\frac{434}{1000},\qquad
 \frac{306}{1000}<\frac{\sqrt6}{8}<\frac{307}{1000}
\]
are certified by squaring their positive endpoints.  They are
intervals for algebraic numbers, not exact rational values.
Theorem~\ref{thm:spectralvalue94} and
Corollary~\ref{cor:messagevalue94}, whose no-message reduction is
Lemma~\ref{lem:postponement96}, give both upper bounds in
\eqref{eq:qutritvalues96}.

For explicit attainment use
\[
 a_1=\diag(3/4,1/4,0),\quad
 a_2=\diag(3/4,0,1/4),\quad w_1=w_2=1/2,
\]
so $\rho=(a_1+a_2)/2$.  With $C_0=\sqrt\rho:\C^3\to\C^3$, set
\[
 \begin{array}{ll}
 V_1=\begin{pmatrix}1/\sqrt2&0&0\\0&1&0\end{pmatrix},&
 D_1=\begin{pmatrix}1/\sqrt2&0&0\\0&1/\sqrt2&0\end{pmatrix},\\[6pt]
 V_2=\begin{pmatrix}1/\sqrt2&0&0\\0&0&1\end{pmatrix},&
 D_2=\begin{pmatrix}1/\sqrt2&0&0\\0&0&1/\sqrt2\end{pmatrix}.
 \end{array}
\]
Every displayed map has domain $\C^3$ and codomain $\C^2$.
The old coefficients are $C_h=V_hC_0$, and directly
\[
 \sum_hV_h^*V_h=I_3,\quad C_h^*C_h=a_h/2,\quad
 \tr D_h^*D_h=1.
\]
Prepare $|D_h\rangle$ independently given the receiver outcome.
On $\C^2\otimes\C^2$ let
$|\psi^-\rangle=(|1,2\rangle-|2,1\rangle)/\sqrt2$.
For $y=z$ declare $C$ on
$|\psi^-\rangle\langle\psi^-|$; for $y\ne z$ its effect is zero.
Use the complement for the alternative.  Each branch filtered swap
has negative eigenvalue $-\sqrt3/16$, with eigenvector $|\psi^-\rangle$.
The payoff coefficient is $\kappa=1/21$, so the total improvement,
with three first labels and two outcomes, is
$3\cdot2\cdot\sqrt3/(21\cdot16)=\sqrt3/56$.
The same instrument is used after every first label.

Without a message keep $C_0$ and use $D_1$ after every first label.
The negative eigenvector is the antisymmetric vector in the first two
old coordinates and the two fresh coordinates; its eigenvalue is
$-\sqrt6/16$.  The same equal-label readout gives improvement
$3\cdot\sqrt6/(21\cdot16)=\sqrt6/112$.
The squared-norm normalizations above, not an unnormalized Bell
vector, determine these factors.  All outcomes are included and no
initial mixture or favorable-history selection has been introduced.
\end{proof}
The same construction with $0<t<1$ uses full-rank devices; its gain is
$t^2(2\sqrt3-\sqrt6)/[16(8-t^2)]$.  At $t=0$ it is zero and both
values equal $1/2$.  A positive gain at this fixed $\rho$ does not
contradict the globally optimized no-message attainments at other
initial spectra.  The retained-register interpretation presupposes
the specified calibrated reset cut.

\section{Equal-prior values with a two-dimensional retained receiver}
\label{sec:equalinitial95}
The prescribed-spectrum formula in Section~\ref{sec:spectralvalue94}
uses the biased prior of the first basis experiment.  Equal priors change
not only the baseline but also the optimal fresh acquisition.  We now
solve the equal-prior problem at every pure initial acquisition when the
old receiver has dimension at most two.  The input dimension is arbitrary,
and allowing a larger fresh reference does not change the answer.
The unrestricted qubit instance consequently admits a complete
initial-spectrum formula, including an optimal protocol without a
receiver message.

Use a weighted exact second-moment family as in
\eqref{eq:basismoments90}.  With probability $1/2$ the device is
$C_y=I/d$; otherwise draw $b$ with law $w_b$ \emph{once}, and use
$E_y^b(t)=(1-t)I/d+tP_y^b$ at both calls.  Here $d\ge2$ and
$0\le t\le1$.  The initial acquisition is the prescribed pure vector
$|C_0\rangle$, with Gram matrix $\rho=C_0^*C_0\in\mathcal D_d$;
its input marginal is $\rho^{\mathsf T}$.  It is not replaced by a
classically resolved initial mixture.  All instruments, discarded calls,
public records, null histories and norm limits follow
Definition~\ref{def:rankclass92}.  Thus the complete retained old system,
not merely one of its subsystems, has dimension at most two.

For $0\le v\le1$ define
\begin{equation}\label{eq:psidef95}
 \begin{split}
 c_d(v)&=d(d-2)+(2d-1)v,\\
 \psi_d(v)&=d-1+
 \frac{(2d-1)v}{2\left(d-1+
               d\sqrt{c_d(v)/(d^2-1)}\right)}.
 \end{split}
\end{equation}
The denominator is positive, including at $d=2,v=0$.
Set
\[
 q=\max\{\lambda_{\max}(\rho),1/2\},\qquad v_\rho=4q(1-q).
\]
The pair $(q,1-q)$ is the nonzero part, with a possible zero entry,
of the least rank-two majorant \eqref{eq:rankmajorant94}.

\begin{theorem}[Equal-prior prescribed-spectrum value]
\label{thm:equalinitial95}
For the once-selected experiment and the fixed pure first acquisition
above, and every $2\le\ell\le d$, the full class
$\mathfrak R_{2,\ell}$ has exact optimum
\begin{equation}\label{eq:equalinitial95}
 P_{2,\ell}^{\rm eq}(\rho)
   =\frac12+\frac{t^2}{2d(d+1)}\psi_d(v_\rho).
\end{equation}
A recorded instrument with at most $d$ outcomes attains this value on
the prescribed initial receiver.  The same instrument may be used after
every first device label; both retained quantum registers may have
dimension two.  The fresh preparation depends only on its recorded
receiver outcome, not on the first device label.  Its Schmidt weights
are specified in Lemma~\ref{lem:ranktwoleaf95} below.

For $t>0$ the value equals the optimized rank-two benchmark precisely
when $\lambda_{\max}(\rho)\le1/2$, and decreases strictly as that
largest eigenvalue increases above $1/2$.  At $t=0$ every value is
$1/2$.  If either retained register has dimension at most one, the
exact value at every fixed $\rho$ is instead
\begin{equation}\label{eq:rankoneinitial95}
 \frac12+\frac{t^2(d-1)}{2d(d+1)}.
\end{equation}
\end{theorem}
The first assertion covers all fresh dimensions at least two while
fixing the old dimension at two.  It makes no assertion that an old
receiver of larger dimension has the same value.  The general-rank
biased formula and the equal-prior globally optimized dimension
hierarchy retain their respective quantifiers.  If $b$ is redrawn
independently at the two calls, the averaged alternative is
$\mathcal M_C\otimes\mathcal M_C$, and all equal-prior values in
this section are $1/2$ instead.

\subsection{A concavity principle for filtering}
For a Hermitian $X$ and a positive matrix $M$ on a finite-dimensional
space, let $\mathcal F_X(M)=\|\sqrt M X\sqrt M\|_1$.
\begin{lemma}[Concavity in the filter Gram]\label{lem:filterconcavity95}
For every $M\succeq0$,
\begin{equation}\label{eq:filterconcavity95}
 \mathcal F_X(M)
  =\max_{-M\preceq Z\preceq M}\tr(XZ).
\end{equation}
In particular $\mathcal F_X$ is continuous, positively homogeneous
and concave on the positive cone.  Consequently, for fixed $A\succeq0$,
\[
 B\longmapsto
 \|(\sqrt A\otimes\sqrt B)(I-dS)(\sqrt A\otimes\sqrt B)\|_1
\]
is concave.  If $A$ is diagonal, replacing $B$ by its diagonal
pinching cannot decrease this last quantity.
\end{lemma}
\begin{proof}
Every Hermitian contraction $F$ gives
$Z=\sqrt M F\sqrt M$ in the interval $[-M,M]$.
Conversely any matrix in this interval vanishes on $\ker M$.
Indeed $M\pm Z\succeq0$ have zero quadratic form on that kernel,
so both annihilate it.  On $\operatorname{supp}M$,
$F=M^{-1/2}ZM^{-1/2}$ is a Hermitian contraction; extend it by zero.
The variational formula for the trace norm and cyclicity of the trace
prove \eqref{eq:filterconcavity95}, without invertibility of $M$.
Convex combinations of feasible optimizers are feasible for the same
combination of their bounds $M$.  This proves concavity.  Homogeneity
is immediate and continuity follows from the square-root formula.

Apply the result to $M=A\otimes B$.  If $U$ is a diagonal unitary,
$U$ commutes with $A$ and $(U\otimes U)S=S(U\otimes U)$.
Thus simultaneous conjugation shows that the value at $UBU^*$ equals
that at $B$.  Averaging diagonal sign unitaries eliminates every
off-diagonal entry, and concavity proves the pinching assertion.
Pinching may increase rank; it is used below only in an
\emph{unrestricted} upper bound.  A rank-two optimizer is subsequently
constructed, rather than inferred from a rank-preserving pinching.
\end{proof}

\subsection{The exact rank-two leaf}
\begin{lemma}[A centered swap with a rank-two old Gram]
\label{lem:ranktwoleaf95}
Let $A\succeq0$ have rank at most two and trace $a>0$.  Let
$x,1-x$ be the eigenvalues of $A/a$ on a two-dimensional support,
padded if necessary, with $1/2\le x\le1$.  For every $\ell\ge2$,
\begin{equation}\label{eq:ranktwoleaf95}
 \max_{B\in\mathcal D_d^{\le\ell}}
 \|(\sqrt A\otimes\sqrt B)(I-dS)(\sqrt A\otimes\sqrt B)\|_1
                 =a\,\psi_d(4x(1-x)).
\end{equation}
The same maximum holds without a rank restriction on $B$.  It has
an optimizer commuting with $A$ and supported on its leading two
coordinates.  For $v=4x(1-x)>0$, put
\begin{equation}\label{eq:freshweights95}
 z=2x-1,\quad e=d-1,\quad
 T=2\psi_d(v)-e,\quad
 w_*=\frac{z(eT-1)}{c_d(v)}.
\end{equation}
The two fresh eigenvalues are $(1+w_*)/2,(1-w_*)/2$ in that
ordered eigenbasis.  Here $0\le w_*\le z<1$.  At rank one, an
aligned pure fresh Gram is optimal and the value is $a(d-1)$.
At $A=0$ the value is zero.
\end{lemma}
\begin{proof}
Divide by $a$ and diagonalize $A$.
Lemma~\ref{lem:filterconcavity95} permits a diagonal $B$ for the
unrestricted maximum.  If $A$ has rank one, the norm is
$(d-1)b_1+\sum_{j>1}b_j\le d-1$, attained by $b_1=1$.
For rank two, write $m=b_1+b_2$.  The part with the fresh coordinate
outside the support of $A$ is positive and has trace $1-m$;
the swap term vanishes there.  The remaining part has norm $m$
times its normalized two-coordinate value, when $m>0$.
The two-coordinate maximum is at least $d-1$: for example $B=A$
gives
\[
 (d-1)\bigl(x^2+(1-x)^2\bigr)+2d x(1-x)
                   =d-1+2x(1-x).
\]
It follows that all fresh mass may be placed in those two coordinates.
This also constructs a candidate of rank at most two and justifies
applying the unrestricted upper at every $\ell\ge2$.

Write the two fresh weights as $(1+w)/2,(1-w)/2$, $-1\le w\le1$.
The vectors $|11\rangle,|22\rangle$ give the negative diagonal
entries $(1-d)x(1+w)/2$ and $(1-d)(1-x)(1-w)/2$.
On $|12\rangle,|21\rangle$ the block is
\[
 \begin{pmatrix}
 x(1-w)/2&-\tfrac d2\sqrt{x(1-x)(1-w^2)}\\
 -\tfrac d2\sqrt{x(1-x)(1-w^2)}&(1-x)(1+w)/2
 \end{pmatrix}.
\]
Its determinant is nonpositive.  Hence the complete trace norm is
\begin{equation}\label{eq:scalarleaf95}
 F_z(w)=\frac12\left[e(1+zw)+
       \sqrt{(z-w)^2+d^2v(1-w^2)}\right],\qquad v=1-z^2.
\end{equation}
This formula includes $w=\pm1$ by continuity.

Here is an algebraic upper certificate, which also evaluates the
maximum.  Set $C=c_d(v)>0$ and
$D(w)=1+(d^2-1)v-2zw+(1-d^2v)w^2$.
Substitution of \eqref{eq:psidef95}--\eqref{eq:freshweights95} gives
\begin{equation}\label{eq:scalarcertificate95}
 (T-ezw)^2-D(w)=C(w-w_*)^2.
\end{equation}
For clarity, its quadratic coefficient is
$e^2z^2-(1-d^2v)=C$, its linear coefficient is
$-2z(eT-1)$, and its constant coefficient equals
$z^2(eT-1)^2/C$.  The last identity follows by squaring
$\sqrt{C/(d^2-1)}$ in the expression for $T$.
Since $T\ge e$ and $z<1$, $T-ezw>0$ throughout $[-1,1]$.
Moreover
\[
 eT-1=d(d-2)+
     \frac{e(2d-1)v}{e+d\sqrt{C/(d^2-1)}}\in[0,C],
\]
which proves $0\le w_*\le z<1$.
Taking square roots in \eqref{eq:scalarcertificate95} proves
$F_z(w)\le(e+T)/2=\psi_d(v)$, with equality at $w_*$.
No squared stationary equation with an unchecked sign is used.
The rank-one and zero cases were treated before division by $C$.
Finally polar isometries preserve these trace norms for all rectangular
coefficient maps having Grams $A,B$.
\end{proof}

\begin{lemma}[Concavity of the rank-two spectral objective]
\label{lem:psiconcavity95}
The function $\psi_d$ is continuous, strictly increasing and concave
on $[0,1]$, with
\[
 \psi_d(0)=d-1,\qquad\psi_d(1)=d-1/2.
\]
The function $f_d(x,1-x)=\psi_d(4x(1-x))$ is symmetric and
strictly concave on the two-coordinate probability simplex.
\end{lemma}
\begin{proof}
In the positive added term of \eqref{eq:psidef95} write
\[
 a=d-1,\quad b=d/\sqrt{d^2-1},\quad c=d(d-2),\quad s=2d-1.
\]
It suffices to examine $h(v)=v/(a+b\sqrt{c+sv})$.
For $y=\sqrt{c+sv}>0$ differentiation gives
\[
 h'(v)=\frac{a+(b/2)(y+c/y)}{(a+by)^2}>0.
\]
The derivative of this last expression with respect to $y$ is
\[
 -\frac{b}{2(a+by)^3}
       \left(3a+\frac{ac}{y^2}+by+\frac{3bc}{y}\right)<0.
\]
Since $y$ increases with $v$, this proves concavity in the interior.
When $d=2$ the formula extends continuously to $v=0$ with finite
right derivative; concavity on the closed interval follows by limits.
The endpoint values follow by substitution.  Finally $4x(1-x)$
is strictly concave and symmetric, and composition with an increasing
concave function gives the claimed strict concavity of $f_d$.
\end{proof}

\subsection{Full protocol optimization and attainment}
\begin{proof}[Proof of Theorem~\ref{thm:equalinitial95}]
Write $\alpha=t^2/[2d^2(d+1)]$.
The two equal-prior payoff blocks are those of
\eqref{eq:equalpayoff91}: $\alpha(I-dS)$ for equal labels and
$-\alpha(I-dS)/(d-1)$ for unequal labels.
Use the dimension-preserving normal form of
Lemma~\ref{lem:ranknormal92}, on the prescribed initial acquisition.
For each first label, its Grams satisfy
\[
 \sum_h A_{yh}=\rho,\quad \operatorname{rank}A_{yh}\le2,
 \quad \tr b_{yh}=1,\quad\operatorname{rank}b_{yh}\le\ell.
\]
On a nonzero branch put $w_{yh}=\tr A_{yh}$ and
$a_{yh}=A_{yh}/w_{yh}$.  These are Gram weights, not conditional
history probabilities.  Optimizing final effects over each second
label gives exactly $\alpha$ times the trace norm of the centered
filtered swap: the $d-1$ opposite-sign blocks cancel their denominator,
and $\tr X_++\tr(-X)_+=\|X\|_1$.
Lemma~\ref{lem:ranktwoleaf95} therefore bounds the gain by
\[
 \alpha\sum_{y,h}w_{yh}f_d(\lambda(a_{yh})).
\]
Apply Lemma~\ref{lem:spectralroof94} with rank two and the
continuous symmetric concave function of
Lemma~\ref{lem:psiconcavity95}.  For every $y$ the corresponding
sum is at most $f_d(q,1-q)=\psi_d(v_\rho)$.
There are $d$ first labels.  This proves the upper in
\eqref{eq:equalinitial95}.  Kraus refinements of mixed receiver maps
and classical refinements of fresh mixtures do not retain an extra
quantum environment.  Countable outcomes and null histories use
Lemma~\ref{lem:countable91}; the bound then passes to tester closure.

For attainment choose the at-most-$d$ commuting atoms of
Lemma~\ref{lem:spectralroof94},
$\rho=\sum_hw_h a_h$, all of spectrum $(q,1-q,0,\ldots,0)$.
Factor $w_h a_h=C_h^*C_h$ with $C_h:\C^d\to\C^2$.
The maps $V_h=C_hC_0^+$ satisfy
$V_hC_0=C_h$ and $\sum_hV_h^*V_h=I$ on the initial Schmidt
support, and extend to a complete recorded instrument on the initial
receiver.  For each $h$, choose the commuting maximizing $b_h$
of Lemma~\ref{lem:ranktwoleaf95} and a coefficient map
$D_h:\C^d\to\C^2$ with $D_h^*D_h=b_h$.
Prepare its normalized vector independently of the old output,
conditional on $h$.  For the centered receiver operator $X_h$,
declare $C$ on its positive spectral subspace when labels agree,
and on its negative spectral subspace when they differ; assign the
kernel arbitrarily.  The complementary effects complete a binary
POVM on every block.  Every leaf bound and the envelope bound is now
attained.  The instrument is common to all first labels and is defined
without dividing by any physical branch probability.

The endpoint and strictness claims follow from
Lemma~\ref{lem:psiconcavity95}.  For the rank-one assertion,
Theorem~\ref{thm:rankspectrum92} gives the upper independently of
$\rho$.  Decompose this fixed Gram into pure spectral atoms and
construct the same support-inverse recorded instrument.  Prepare an
aligned pure fresh input for each atom.  Its centered norm is $d-1$,
so the positive-part decisions attain that upper with both receivers
of dimension one.  This proves the assertion whenever either of the
allowed dimensions is one.  The one-call channel identity and
independent-redraw conclusion follow from the first moment, exactly as
in Theorem~\ref{thm:equalprior91}.
\end{proof}

\subsection{A complete qubit formula and spectral inference}
\begin{corollary}[Every pure initial qubit acquisition]
\label{cor:qubitinitial95}
For $d=2$ and any $\rho\in\mathcal D_2$, put
$u=\sqrt{\det\rho}$.  With both retained registers allowed dimension
two, the exact equal-prior reset value is
\begin{equation}\label{eq:qubitinitial95}
 P_{2,2}^{\rm eq}(\rho)
  =\frac12+\frac{t^2}{12}
                     \left(1+\frac{6u^2}{1+4u}\right).
\end{equation}
After identifying the initial Schmidt support with a subspace of $\C^2$,
no nontrivial recorded receiver instrument or receiver message is required.
In an eigenbasis of $\rho$, write its eigenvalues as
$(1+z)/2,(1-z)/2$, $0\le z\le1$.  An optimal fresh Gram has weights
\begin{equation}\label{eq:qubitfresh95}
 \frac12\left(1+\frac{z}{1+2\sqrt{1-z^2}}\right),\qquad
 \frac12\left(1-\frac{z}{1+2\sqrt{1-z^2}}\right).
\end{equation}
The final joint spectral decisions are as in the theorem.
All receiver feedback, public randomization and intermediate receiver
processing are included in the upper bound.
\end{corollary}
\begin{proof}
Here the initial rank is at most two, so its rank-two majorant is its
own spectrum.  In the preceding attainment the single atom $\rho$
suffices: keep the first receiver unchanged on that Schmidt support.  In
\eqref{eq:psidef95}, $v_\rho=4\det\rho=4u^2$ and
$\psi_2(v)=1+3v/[2(1+2\sqrt v)]$.
This gives \eqref{eq:qubitinitial95}.
For $v>0$, \eqref{eq:freshweights95} simplifies to
$w_*=z/(1+2\sqrt v)$; its endpoint limit gives the stated pure
fresh Gram at $z=1$.
\end{proof}
Thus a pure initial Gram gives $1/2+t^2/12$ and the maximally mixed
Gram gives $1/2+t^2/8$.  At a strictly intermediate spectrum the
maximizing fresh spectrum is generally neither the initial spectrum
nor the uniform spectrum.  This differs from the biased rank-two
leaf, whose optimal fresh weights are equal.

\begin{corollary}[Exact loss and a calibrated spectral bound]
\label{cor:equalspectralloss95}
In Theorem~\ref{thm:equalinitial95}, the exact loss from the optimized
rank-two value is
\begin{equation}\label{eq:equalspectralloss95}
 \frac{t^2}{2d(d+1)}
                \left(d-\frac12-\psi_d(v_\rho)\right).
\end{equation}
Suppose $t>0$ and an implemented score $p$ differs by at most
$\eta\ge0$ from the score of a nominal strategy with this fixed
initial acquisition and the stated reset dimensions.  Put
\[
 R=\frac{2d(d+1)}{t^2}\left(p-\eta-\frac12\right).
\]
If $d-1<R\le d-1/2$, let $v_R\in(0,1]$ be the unique solution
$\psi_d(v_R)=R$.  Then
\begin{equation}\label{eq:equalcapinference95}
 \lambda_{\max}(\rho)\le\frac{1+\sqrt{1-v_R}}2.
\end{equation}
If $R\le d-1$ this reasoning gives only the trivial bound one.
If $R>d-1/2$, the nominal model and the supplied error bound cannot
both describe that score.
\end{corollary}
\begin{proof}
Subtract \eqref{eq:equalinitial95} from
\eqref{eq:rankequal92} at rank two.  For the second assertion,
$p-\eta\le P_{2,\ell}^{\rm eq}(\rho)$ gives
$R\le\psi_d(v_\rho)$.  Strict monotonicity gives
$v_\rho\ge v_R$, and solving $4q(1-q)\ge v_R$ for $q\ge1/2$
proves \eqref{eq:equalcapinference95}.  The two remaining cases follow
from the endpoint range of $\psi_d$.
\end{proof}
The error hypothesis concerns a specified implementation and nominal
experiment, with fixed priors and latent-index law.  A uniform
unhalved diamond error $\delta$ on each complete device channel
contributes at most $\delta$ to a two-call decision score; a certified
hard-path preparation error $\varepsilon$ contributes at most
$\varepsilon/2$.  Their sum may therefore be used for $\eta$, as
in Corollary~\ref{cor:gamestability91}.  A fidelity estimate, a sampled
matrix inequality or an uncalibrated score is not substituted for this
hypothesis.  The bound infers a property of the prescribed initial Gram,
not its eigenbasis or the dimension of an unrestricted coherent machine.

\section{Comparison of operational models}
\label{sec:comparison96}\label{sec:comparison84}
The feasible set, not merely the word ``memory'', determines a decision
problem.  We record four comparisons relevant to the general and
spectral results above.

\paragraph{Constrained-separability testers.}
Ohst, Zhang, Nguyen, Pl\'avala and Quintino
\cite[Sections~5.3--5.4 and~6.1--6.2]{OZNPQ89} formulate auxiliary
memory restrictions by positive separable decompositions with affine
conditions on their factors.  Their Definition~23 and Theorem~24
characterize complete measure-and-reprepare adaptive testers and their
constrained-separability formulation.  Our class $\mathfrak C$ uses
that complete-readout boundary.  A reset tester retaining an old
receiver instead need not have separable final complete-call effects:
the later joint readout is unrestricted, while coherent import into
the new source is prohibited.  Our fixed-initial objective constrains
the common Gram barycenter and evaluates it explicitly for the
shared-basis payoff.  The general tester framework is an antecedent,
not a claim superseded by a first construction of memory hierarchies.

\paragraph{Autonomous coherent memory.}
Zonnios and Binder \cite[Proposition~2, Theorem~1 and Corollary~1]{ZB90}
optimize over a repeated stationary instrument carrying a coherent
register between steps and retaining the classical record.  Their
finite-time completeness compiles a general adaptive tester into an
autonomous machine with sufficiently large memory and a counter.
Here the source cannot receive the old coherent register, public
controls may depend on the step, and the two specified receiver cuts
are not an all-time workspace dimension.  Saturation at
$\operatorname{rank}\rho$ in our biased experiment concerns the
exact value of that family with its first acquisition fixed.  It is
not their arbitrary-horizon compilation theorem.

\paragraph{Majorization and ensemble transformations.}
Nielsen \cite{Nielsen94} characterizes deterministic pure-state
transformations by majorization.  Jonathan and Plenio
\cite[Theorem~1]{JonathanPlenio94} treat ensembles of pure-state
outputs through averaged Schmidt constraints.  These are direct
antecedents for the spectral barycenter step.  Our feasible
instruments act only on the already prepared receiver, have the common
sum $\sum_hA_{yh}=\rho$ after each first device label, and optimize
subsequent measurement decisions.  Their transformation targets are
not the payoff matrices \eqref{eq:gameblocks91}.  The least-majorant
argument, finite convexity and support inverse are used here with
credit; the specific contribution is their combination with the exact
rank-constrained filtered-swap leaf and its message comparator.

\paragraph{Assistance and concave roofs.}
Entanglement assistance \cite{Assistance94} and general roof methods
\cite{Uhlmann91} optimize averages over decompositions of a fixed
state.  The present all-density and pure-atom roofs have such a
mathematical form, but their atoms, rank cuts and continuation reward
come from the reset normal form.  The uniform message criterion asks
whether these continuation functions are concave; the binary witness
changes one receiver-to-source classical permission, not an
entanglement-conversion task.  A universal majorant is a condition on
all admissible atoms and cannot be replaced by finite sampling.

Finally, a measurement channel with classical output is not a L\"uders
instrument exposing a post-measurement state \cite{EQ90}.
The local geometry in the supplement, the general variational theory,
and the closed shared-basis formulas remain separate theorem families.
No priority for majorization, concave envelopes, swap tests, or
memory-constrained discrimination as general mechanisms is claimed.

\begin{thebibliography}{99}
\bibitem{SZ74} M. Sedl\'ak and M. Ziman,
Optimal single-shot strategies for discrimination of quantum measurements,
\emph{Phys. Rev. A} \textbf{90} (2014), 052312;
arXiv:1408.0934; doi:10.1103/PhysRevA.90.052312.
Section VI, Example 4, equation (37).
\bibitem{OZNPQ89} T.-A.~Ohst, S.~Zhang, H.~C.~Nguyen,
M.~Pl\'avala and M.~T.~Quintino,
Characterising memory in quantum channel discrimination via
constrained separability problems,
\emph{Quantum} \textbf{10} (2026), 1988;
doi:10.22331/q-2026-01-28-1988; arXiv:2411.08110v2.
\bibitem{ZB90} M.~Zonnios and F.~C.~Binder,
Distinguishing quantum processes with bounded coherent memory,
Preprint, 2026. arXiv:2606.19511v1.
\bibitem{EQ90} C.~Eid and M.~T.~Quintino,
Post-measurement states are (very) useful for measurement discrimination,
arXiv:2602.12258v2 (2026).
\bibitem{GAE90} D.~Gross, K.~Audenaert and J.~Eisert,
Evenly distributed unitaries: on the structure of unitary designs,
arXiv:quant-ph/0611002 (2006).
\bibitem{ZHS90} M.~Ziman, T.~Heinosaari and M.~Sedl\'ak,
Unambiguous comparison of quantum measurements,
arXiv:0905.4445 (2009).
\bibitem{Uhlmann91} A.~Uhlmann, Roofs and convexity,
Entropy \textbf{12} (2010), 1799--1832.
\bibitem{BV91} S.~Boyd and L.~Vandenberghe, \emph{Convex Optimization},
Cambridge University Press, Cambridge, 2004.
\bibitem{FvdG93} C. A. Fuchs and J. van de Graaf,
Cryptographic distinguishability measures for quantum-mechanical states,
\emph{IEEE Trans. Inform. Theory} \textbf{45} (1999), 1216--1227;
arXiv:quant-ph/9712042, Theorem~1.
\bibitem{Nielsen94} M. A. Nielsen,
Conditions for a class of entanglement transformations,
\emph{Phys. Rev. Lett.} \textbf{83} (1999), 436--439;
arXiv:quant-ph/9811053, Theorem~1.
\bibitem{JonathanPlenio94} D. Jonathan and M. B. Plenio,
Minimal conditions for local pure-state entanglement manipulation,
\emph{Phys. Rev. Lett.} \textbf{83} (1999), 1455--1458;
erratum, \textbf{84} (2000), 4781.
\bibitem{Assistance94} D. P. DiVincenzo, C. A. Fuchs, H. Mabuchi,
J. A. Smolin, A. Thapliyal and A. Uhlmann,
Entanglement of assistance, arXiv:quant-ph/9803033 (1998).
\end{thebibliography}

\end{document}
